\documentclass[11pt, reqno]{amsart}

\usepackage[T1]{fontenc}
\usepackage[utf8]{inputenc}
\usepackage[margin=2.4cm]{geometry}
\usepackage{amsmath, amssymb, amsthm}
\usepackage{booktabs}
\usepackage{microtype}
\usepackage[colorlinks=true, citecolor=blue, linkcolor=blue, urlcolor=blue]{hyperref}

\newtheorem{theorem}{Theorem}
\newtheorem{corollary}[theorem]{Corollary}
\newtheorem{lemma}{Lemma}[section]
\newtheorem{proposition}[lemma]{Proposition}

\theoremstyle{remark}
\newtheorem{remark}[lemma]{Remark}

\numberwithin{equation}{section}

\DeclareMathOperator{\re}{Re}
\DeclareMathOperator{\im}{Im}

\newcommand{\Harm}{\mathsf{H}}
\newcommand{\R}{\mathbb{R}}
\newcommand{\C}{\mathbb{C}}
\newcommand{\cH}{\mathcal{H}}
\newcommand{\cR}{\mathcal{R}}
\newcommand{\cG}{\mathcal{G}}
\newcommand{\dd}{\,\mathrm{d}}

\title[Asymptotic bounds on Maynard's sieve functional]{%
Asymptotic lower bounds on Maynard's sieve functional\\
and microscale uniqueness}

\author{Fabiano Fiorin}
\email{fabianofiorin06@gmail.com}

\subjclass[2020]{Primary 11N05, 11N36; Secondary 65G20}
\keywords{Bounded gaps between primes, Maynard--Tao sieve, symmetric power sums,
Laplace transforms, rigorous interval arithmetic}

\begin{document}

\begin{abstract}
We establish a new asymptotic lower bound on Maynard's sieve variational functional
$M_k := \sup_F R_k(F)$ ($R_k(F) := k J(F)/I(F)$) as $k \to \infty$, improving Lewulis's
bound $M_k \ge \log k - 0.273572 + o(1)$ to
\[
M_k \ge \log k - 0.260989 + o(1).
\]
This gain arises from combining two mechanisms: (i) an $O(k)$ shift of the reciprocal
scale $a = k(\log k + \alpha)$ at $\alpha = -523/2000$ that aligns the simplex
truncation $P_1 \le 1$ against an asymmetric $1$-stable limiting density, and
(ii) the complete weight-graded hierarchy $\Lambda_9$ of symmetric power-sum products
$\prod_{r \in \lambda} (\sigma^r P_r)$ over all $30$ integer partitions $\lambda$ of
weight at most $9$ with parts at least $2$ ($N = 250$ polynomial parameters).
Restricting to the Bombieri--Vinogradov level ($\theta = 1/2$) and conditionally
assuming the Elliott--Halberstam conjecture ($\theta = 1$), this improves the respective
leading constants for $H_m := \liminf_{n\to\infty}(p_{n+m}-p_n)$ to
$H_m \le (5.1928517 + o(1))\,m\,e^{4m}$ and $H_m \le (2.5964259 + o(1))\,m\,e^{2m}$
(while unconditionally as $m \to \infty$, equidistribution beyond $\theta = 1/2$ yields
smaller exponential growth rates). Conversely, we prove a Microscale Uniqueness Theorem:
among regular reciprocal-type microscale kernels $g(u)$ normalized to $m_0(g)=1$ and
$g(u) \sim u^{-1}$ as $u \to \infty$, the Cauchy kernel $h(u) = (1+u)^{-1}$ is uniquely optimal
within the product-kernel and radial-polynomial class, with every normalized kernel $g \neq h$
(a.e.) incurring the strict penalty $-\int_0^\infty (1+u)(g(u) - h(u))^2 \dd u < 0$ (so among
unnormalized kernels, $h$ is unique up to multiplicative scaling and positive argument dilation).
All bounds are proved
analytically and deterministically certified in
$768$-bit Arb ball arithmetic via an independent Bernstein-ellipse Gauss--Legendre verifier.
\end{abstract}

\maketitle

\section{Introduction and Main Results}

Let $p_n$ denote the $n$-th prime and set $H_m := \liminf_{n\to\infty} (p_{n+m} - p_n)$
for each integer $m \ge 1$. Following Zhang \cite{Zhang} and Goldston--Pintz--Y{\i}ld{\i}r{\i}m
\cite{GPY}, Maynard \cite{Maynard} and Tao reduced the detection of bounded prime gaps in
admissible $k$-tuples to maximizing the symmetric Rayleigh quotient
$R_k(F) := k J(F)/I(F)$ on the $k$-simplex
$\cR_k := \{(t_1,\dots,t_k) \in [0,\infty)^k : t_1 + \cdots + t_k \le 1\}$, where
\begin{align*}
I(F) &:= \int_{\cR_k} F(t_1,\dots,t_k)^2 \dd t_1\cdots \dd t_k, \\
J(F) &:= \int_{\cR_{k-1}} \left(\int_0^{1-\sum_{i=1}^{k-1} t_i}
F(t_1,\dots,t_k)\dd t_k \right)^{\!2} \dd t_1\cdots \dd t_{k-1},
\end{align*}
and $M_k := \sup_{F} R_k(F)$ ranges over locally square-integrable symmetric functions
with $I(F) > 0$. By \cite[Theorems~22 and~26]{PM8b} (at $\varepsilon = 0$), if primes are equidistributed in
arithmetic progressions at level $\vartheta \in (0, 1)$, then $M_k > 2m/\vartheta$ implies
$H_m \le H(k)$, the minimal diameter of an admissible $k$-tuple; Bombieri--Vinogradov
provides $\vartheta = 1/2$ unconditionally ($M_k > 4m$), while the Elliott--Halberstam conjecture
corresponds to $\theta = 1$ (specifically, $\mathrm{EH}[\vartheta]$ holding for every $\vartheta \in (0, 1)$,
yielding $H_m \le H(k)$ whenever $M_k > 2m$).

Asymptotically as $k \to \infty$, Maynard \cite{Maynard} obtained
$M_k \ge \log k - 2\log\log k - 2 + o(1)$, and Polymath8b \cite[Corollary~37]{PM8b}
established $M_k \ge \log k - O(1)$ together with the upper barrier
\begin{equation}\label{eq:pm8b-barrier}
M_k \le \frac{k}{k-1}\log k \qquad (k \ge 2),
\end{equation}
forcing any lower bound $M_k \ge \log k - c^* + o(1)$ to satisfy $c^* \ge 0$. Combining
equidistribution beyond $\theta = 1/2$ yields strictly smaller unconditional exponential
rates for $H_m$ as $m \to \infty$: Polymath8b \cite{PM8b} proved $H_m \ll m\exp\bigl((4-\tfrac{28}{157})m\bigr)$, and Stadlmann \cite{Stadlmann2025}
improved this via smooth-moduli equidistribution ($\theta = \tfrac{1}{2}+\tfrac{1}{40}-\varepsilon$)
to $H_m = O(\exp(3.8075 m))$ (and later $H_1 \le 240$ \cite{Stadlmann2026}).

Recently, Lewulis \cite{Lewulis} analyzed the exact Laplace representation of $I(F)$ and
$J(F)$ for the radial reciprocal ansatz
$F(t_1,\dots,t_k) = \widetilde{Q}((\log k)(1-P_1)) \prod_{i=1}^k (1 + a t_i)^{-1}$
($P_r := \sum_{i=1}^k t_i^r$) at $a = k\log k$, optimizing over polynomials $\widetilde{Q}$
of degree at most $20$ to prove \cite[Theorem~3]{Lewulis}:
\begin{equation}\label{eq:lewulis-bound}
M_k \ge \log k - 0.273572 + o(1) \qquad (k \to \infty).
\end{equation}

In this paper we both characterize the microscale barrier in~\eqref{eq:lewulis-bound} and
overcome its macroscopic radial limit. Throughout, $\gamma = 0.57721566\dots$ denotes the
Euler--Mascheroni constant, $L := \log k$, and $\Harm_n := \sum_{j=1}^n j^{-1}$ ($\Harm_0 := 0$)
is the $n$-th harmonic number.

\begin{theorem}[Microscale Uniqueness]\label{thm:uniqueness}
Let $h(u) := (1+u)^{-1}$ for $u \ge 0$, and let $g \colon [0,\infty) \to \R$ be a
measurable function satisfying:
\begin{enumerate}
\item[(H1)] $g \in L^2(0,\infty) \cap L^4(0,\infty)$ and $m_0(g) := \int_0^\infty g(u)^2 \dd u > 0$;
\item[(H2)] $g(u) = c u^{-1} + O(u^{-1-\delta})$ as $u \to \infty$ for some $c \in \R \setminus \{0\}$ and $\delta > 0$.
\end{enumerate}
Set $\tilde{g} := c^{-1} g$, $\lambda := m_0(\tilde{g})^{-1} = c^2 / m_0(g) > 0$, and
$g^*(u) := \lambda \tilde{g}(\lambda u) = \frac{c}{m_0(g)} g\bigl(\frac{c^2}{m_0(g)} u\bigr)$
(note that (H1)--(H2) already imply $\int_0^\infty u (g^*(u) - h(u))^2 \dd u < \infty$). Let $\widetilde{Q}$ be any non-zero real polynomial, set
$\sigma := \log k + \alpha$, $a' := k\sigma$, $a := \lambda k\sigma = \lambda a'$, and define
\begin{equation}\label{eq:Fg-def}
F_g(t_1,\dots,t_k) := \widetilde{Q}\bigl(\sigma(1 - P_1)\bigr) \prod_{i=1}^k g(a t_i).
\end{equation}
Then $m_0(g^*) = 1$, $g^*(u) = u^{-1} + O(u^{-1-\delta})$ as $u \to \infty$, and by scale
invariance $R_k(F_g)$ at scale $a = \lambda k \sigma$ equals $R_k(F_{g^*})$ at scale $a'$.
Moreover, for every $\alpha \in \R$, the limit
$\eta(g^*; \alpha, \widetilde{Q}) := \lim_{k\to\infty}(R_k(F_{g^*}) - \log k)$ exists, is finite,
and satisfies the pointwise shift identity
$\eta(g^*; \alpha, \widetilde{Q}) = \bigl(\kappa(g^*) - \gamma\bigr) + \eta(h; \alpha + c_1(g^*) - \gamma, \widetilde{Q})$
for $c_1(g^*) := \gamma - \int_0^\infty u(g^*(u)^2 - h(u)^2)\dd u \in \R$. Since
$\eta(h; 0, \widetilde{Q}) \in \R$ and~\eqref{eq:pm8b-barrier} gives
$\sup_{\alpha'\in\R}\eta(h; \alpha', \widetilde{Q}) \le 0$ (so all suprema are finite), taking
$\sup_{\alpha\in\R}$ gives
\begin{equation}\label{eq:shift-covariance}
\sup_{\alpha \in \R} \eta(g^*; \alpha, \widetilde{Q})
= \bigl(\kappa(g^*) - \gamma\bigr) + \sup_{\alpha' \in \R} \eta(h; \alpha', \widetilde{Q}),
\end{equation}
where
\begin{equation}\label{eq:kappa-penalty}
\kappa(g^*) - \gamma
= -\int_0^\infty (1+u)\bigl(g^*(u) - h(u)\bigr)^2 \dd u \le 0,
\end{equation}
with equality if and only if $g^*(u) = h(u)$ almost everywhere on $(0,\infty)$, i.e.,
$g(u) = C/(1 + \theta u)$ almost everywhere for some constants $C \in \R \setminus \{0\}$
and $\theta > 0$ (uniqueness up to multiplicative scaling and positive argument dilation).
\end{theorem}

Theorem~\ref{thm:uniqueness} proves that no asymptotic improvement beyond~\eqref{eq:lewulis-bound}
is available by varying the reciprocal microscale factor $g(u)$ inside the radial class.
Gains must instead come from the macroscopic polynomial structure on $\cR_k$: shifting $a = k(\log k + \alpha)$
by an $O(k)$ offset $\alpha < 0$ and adjoining higher-order symmetric power sums $P_r = \sum_{i=1}^k t_i^r$.

Let $\Lambda_9$ denote the set of all $30$ integer partitions $\lambda = (\lambda_1, \dots, \lambda_{\ell(\lambda)})$
with $\lambda_1 \ge \cdots \ge \lambda_{\ell(\lambda)} \ge 2$ and $|\lambda| := \sum_{j=1}^{\ell(\lambda)} \lambda_j \le 9$,
grouped by weight $|\lambda| \in \{0, 2, \dots, 9\}$ (including the empty partition $()$ of length $0$):
\begin{equation}\label{eq:Lambda9-list}
\begin{aligned}
\Lambda_9 = \bigl\{ &(),\ (2),\ (3),\ (4),\ (2,2),\ (5),\ (3,2),\ (6),\ (4,2),\ (3,3),\\
&(2,2,2),\ (7),\ (5,2),\ (4,3),\ (3,2,2),\ (8),\ (6,2),\ (5,3),\\
&(4,4),\ (4,2,2),\ (3,3,2),\ (2,2,2,2),\ (9),\ (7,2),\ (6,3),\\
&(5,4),\ (5,2,2),\ (4,3,2),\ (3,3,3),\ (3,2,2,2)\bigr\}.
\end{aligned}
\end{equation}

\begin{theorem}[Certified Asymptotic Bound on $M_k$]\label{thm:main}
Let $L := \log k$, $\sigma := L + \alpha$, and $a := k\sigma$. Taking the symmetric
power-sum ansatz over~\eqref{eq:Lambda9-list},
\begin{equation}\label{eq:ansatz-C}
F(t_1,\dots,t_k) = \left[ \sum_{\lambda \in \Lambda_9}
\widetilde{Q}_\lambda\bigl(\sigma(1-P_1)\bigr)
\prod_{j=1}^{\ell(\lambda)} \bigl(\sigma^{\lambda_j} P_{\lambda_j}\bigr) \right]
\prod_{i=1}^k \frac{1}{1 + a t_i}
\end{equation}
at exact rational shift $\alpha = -523/2000 = -0.2615$ with the certified polynomial vector
$c \in \mathbb{Q}^{250}$ having weight-graded degrees $20, 14, 12, 10, 8, 8, 6, 6, 5$ for
$|\lambda| \in \{0, 2, \dots, 9\}$, respectively ($30$ blocks, $N = 250$ parameters), let
$\eta_* := \frac{c^\top D_{\mathrm{net}}(\alpha) c}{c^\top I(\alpha) c}$ be its certified asymptotic
Rayleigh limit (Proposition~\ref{prop:UL-reduction}, \eqref{eq:quotient-limit}, and~\eqref{eq:eta-cert})
and set $c^* := -\eta_*$. Then unconditionally as $k \to \infty$:
\begin{equation}\label{eq:main-bound}
M_k \ge \log k - c^* + o(1), \qquad c^* \le 0.26098863656469156542.
\end{equation}
More precisely, $768$-bit Arb certification establishes the sharper bound
$c^* \le 0.2609886365646915654192$.
\end{theorem}

Theorem~\ref{thm:main} improves Lewulis's constant~\eqref{eq:lewulis-bound} by more than
$0.012583$. To isolate the two mechanisms behind this gain
in light of Theorem~\ref{thm:uniqueness}, optimizing $\alpha$ alone on Lewulis's radial
degree-$20$ ansatz ($\Lambda_0$ at $\alpha = -3367/10000$, Entry~A in the companion certificate)
lowers $0.273572$ to $c_{\mathrm{rad}}^* \le 0.2654546851$ (a shift gain of $\approx 0.008117$),
while adjoining the non-radial $\Lambda_9$ power sums contributes the remaining
$\approx 0.004466$ reduction down to~\eqref{eq:main-bound}. The consequences for $H_m$ follow.

\begin{corollary}[Asymptotic Prime-Gap Constants]\label{cor:Hm}
From Theorem~\ref{thm:main} and \cite[Theorems~22 and~26]{PM8b}:
\begin{enumerate}
\item[(i)] \textbf{Pure Bombieri--Vinogradov regime ($\theta = 1/2$, not using $\theta > 1/2$):}
As $m \to \infty$,
\[
H_m \le \bigl(4 e^{c^*} + o(1)\bigr)\,m\,e^{4m} \le \bigl(5.1928517 + o(1)\bigr)\,m\,e^{4m}.
\]
\item[(ii)] \textbf{Elliott--Halberstam regime ($\theta = 1$, conditional):} Assuming the
Elliott--Halberstam conjecture, as $m \to \infty$,
\[
H_m \le \bigl(2 e^{c^*} + o(1)\bigr)\,m\,e^{2m} \le \bigl(2.5964259 + o(1)\bigr)\,m\,e^{2m}.
\]
\end{enumerate}
Unconditionally as $m \to \infty$, prime equidistribution beyond $\theta = 1/2$ \cite{PM8b, Stadlmann2025}
yields strictly smaller exponential growth rates down to $H_m = O(\exp(3.8075 m))$ \cite{Stadlmann2025}.
\end{corollary}

\begin{proof}
Fix $\theta \in \{1/2, 1\}$. For integers $k \ge 3$, define
$\varepsilon(k) := \max(0,\, \log k - c^* - M_k)$; by~\eqref{eq:main-bound} of Theorem~\ref{thm:main},
$\varepsilon(k) \to 0$ as $k \to \infty$. Put $K_m := \lceil \exp(2m/\theta + c^*) \rceil$
and $\delta_m := \sup_{k \ge K_m} \varepsilon(k) + m^{-1} \to 0^+$ as $m \to \infty$.
Choosing $k_m := \lceil \exp(2m/\theta + c^* + \delta_m) \rceil \ge K_m$ gives
$\varepsilon(k_m) \le \delta_m - m^{-1}$, whence
\[
M_{k_m} \ge \log k_m - c^* - \varepsilon(k_m)
\ge \left(\frac{2m}{\theta} + c^* + \delta_m\right) - c^* - \bigl(\delta_m - m^{-1}\bigr)
= \frac{2m}{\theta} + \frac{1}{m} > \frac{2m}{\theta}.
\]
By \cite[Theorems~22 and~26]{PM8b} (at $\varepsilon = 0$; when $\theta = 1$, assuming
$\mathrm{EH}[\vartheta]$ for every $\vartheta \in (0, 1)$ and setting
$\vartheta_m := \frac{2m}{2m + 1/(2m)} \in (0,1)$ gives
$M_{k_m} \ge 2m + \frac{1}{m} > 2m + \frac{1}{2m} = \frac{2m}{\vartheta_m}$), this implies
$H_m \le H(k_m)$. Taking the admissible tuple of the first $k_m$ primes strictly greater than
$k_m$ gives $H(k_m) \le (1+o(1)) k_m \log k_m$ by the Prime Number Theorem. Substituting
$k_m = (1+o(1)) e^{c^*} e^{2m/\theta}$ and $\log k_m = (2/\theta + o(1)) m$ yields
$H_m \le (\frac{2}{\theta} e^{c^*} + o(1)) m e^{2m/\theta}$. Evaluating $4 e^{c^*}$ and
$2 e^{c^*}$ for $c^* = -\eta_*$ (outer-rounded to $20$ decimal digits from the $768$-bit Arb ball)
gives
\begin{align*}
4 e^{c^*} &\in [5.1928516527656657000 \pm 3.93\times 10^{-20}] < 5.1928517, \\
2 e^{c^*} &\in [2.5964258263828328500 \pm 1.97\times 10^{-20}] < 2.5964259. \qedhere
\end{align*}
\end{proof}

\begin{table}[ht]
\centering
\caption{Comparison of asymptotic bounds on $\log k - M_k$ and implied
leading sieve constants under pure Bombieri--Vinogradov ($\theta = 1/2$) and Elliott--Halberstam
($\mathrm{EH}$); the PM8b row gives the sieve-method upper barrier~\eqref{eq:pm8b-barrier} on $4e^{c^*}$ and $2e^{c^*}$.}\label{tab:summary}
\small
\setlength{\tabcolsep}{4pt}
\begin{tabular}{@{}lccc@{}}
\toprule
Method / Result & Bound on $\log k - M_k$ & Pure BV ($4e^{c^*}$) & EH ($2e^{c^*}$) \\
\midrule
Maynard \cite{Maynard} & $\le 2\log\log k + 2 + o(1)$ & div.\ ($O(m^3 e^{4m})$) & div.\ ($O(m^3 e^{2m})$) \\
Lewulis \cite{Lewulis} ($\alpha = 0$, $\deg 20$) & $\le 0.273572 + o(1)$ & $5.2586081$ & $2.6293041$ \\
\textbf{Theorem~\ref{thm:main}} ($\alpha = -\tfrac{523}{2000}$, $\Lambda_9$, $N = 250$) & $\mathbf{\le 0.2609886366} + o(1)$ & $\mathbf{5.1928517}$ & $\mathbf{2.5964259}$ \\
\midrule
PM8b sieve barrier \cite{PM8b} & $\ge -\tfrac{\log k}{k-1}$ & $\ge 4$ & $\ge 2$ \\
\bottomrule
\end{tabular}
\end{table}

Table~\ref{tab:summary} places Theorem~\ref{thm:main} and Corollary~\ref{cor:Hm} alongside prior
bounds. Section~\ref{sec:uniqueness} proves Theorem~\ref{thm:uniqueness}; Section~\ref{sec:laplace}
derives the exact finite-$k$ set-partition Laplace representation for general $\Lambda_W$;
Section~\ref{sec:asymptotic} passes to the $k \to \infty$ limit and reduces it to exact
rational quadratic forms in $1$-dimensional moments; Section~\ref{sec:cert} presents the
$768$-bit Arb certification and independent finite-$k$ verification; and Section~\ref{sec:repro}
records reproducibility metadata.

\section{Proof of the Microscale Uniqueness Theorem}\label{sec:uniqueness}

Throughout the paper, $\log z$ denotes the principal branch of the complex logarithm on
$\C \setminus (-\infty, 0]$. For $\alpha \in \R$ and $z = 1 + i\tau$ ($\tau \in \R$), define
\begin{equation}\label{eq:Phi-nu-def}
\Phi_\alpha(z) := z(\gamma + \alpha + \log z),
\qquad
\nu_\alpha(x) := \frac{1}{2\pi}\int_{-\infty}^\infty e^{\Phi_\alpha(1+i\tau) - (1+i\tau)x}\dd\tau.
\end{equation}
Since $\re\Phi_\alpha(1+i\tau) = \gamma + \alpha + \frac{1}{2}\log(1+\tau^2) - \tau\arctan\tau$,
we have
\begin{equation}\label{eq:Phi-abs}
\bigl|e^{\Phi_\alpha(1+i\tau)}\bigr|
= e^{\gamma+\alpha} \sqrt{1+\tau^2}\, e^{-\tau\arctan\tau},
\end{equation}
so $\nu_\alpha(x)$ is continuous on $\R$ and decays at least as $O_\alpha(e^{-x})$.

\begin{lemma}[Positivity and Analyticity of $\nu_\alpha$]\label{lem:nu-pos}
For every $\alpha \in \R$, $\nu_\alpha$ extends holomorphically to the strip
$\{x \in \C : |\im x| < \pi/2\}$, satisfies $\nu_\alpha(x) \ge 0$ for all $x \in \R$,
and is not identically zero on $(0,\infty)$. Hence for any non-zero polynomial $\widetilde{Q}$,
$\int_0^\infty \widetilde{Q}(x)^2 \nu_\alpha(x)\dd x > 0$.
\end{lemma}

\begin{proof}
If $|\im x| \le \frac{\pi}{2} - \eta$ ($\eta > 0$), then $\arctan|\tau| \ge \frac{\pi}{2} - \frac{1}{|\tau|}$
gives $\re(-(1+i\tau)x) - |\tau|\arctan|\tau| \le -\re x - \eta|\tau| + 1$. Thus the integrand
in~\eqref{eq:Phi-nu-def} is bounded locally uniformly by $O_{\alpha,\eta}(e^{-\re x}(1+|\tau|)e^{-\eta|\tau|})$,
proving holomorphicity on $\{|\im x| < \pi/2\}$ and real-analyticity on $\R$.

By Lemma~\ref{lem:W0-props} below (whose proof depends solely on the definition of $W_0$ and
is independent of Section~\ref{sec:uniqueness}), $W_0(s; a)$ for $a = k\sigma$ satisfies
$W_0(\sigma(1+i\tau); k\sigma)/W_0(\sigma; k\sigma) = W_0((1+i\tau)/k; 1)/W_0(1/k; 1)$.
Let $C_k(X; a) \ge 0$ be the $k$-fold convolution density of $t \mapsto (1+at)^{-2}$ ($t \ge 0$).
Since $t \mapsto (1+at)^{-2} \in L^1 \cap L^2$, Plancherel gives $W_0(\sigma+i\omega; a) \in L^2(\R_\omega)$,
so $W_0(\sigma+i\omega; a)^k \in L^1(\R_\omega)$ for $k \ge 2$ and Laplace inversion at
$X = 1 - x/\sigma$ ($s = \sigma(1+i\tau)$) reads
\[
\frac{C_k(1 - x/\sigma; k\sigma)}{\sigma e^{\sigma - (\gamma+\alpha)} W_0(\sigma; k\sigma)^k}
= \frac{1}{2\pi}\int_{-\infty}^\infty e^{i\sigma\tau - (1+i\tau)x + (\gamma+\alpha)}
\left(\frac{W_0((1+i\tau)/k; 1)}{W_0(1/k; 1)}\right)^{\!k}\dd\tau.
\]
Splitting $\tau \in \R$ at $T := 100 L$ and $4k$ for large $k$: on $|\tau| \le T$,
Lemma~\ref{lem:W0-props}(iv) expresses the integrand as
$e^{\Phi_\alpha(1+i\tau)-(1+i\tau)x}(1 + O(L^4/k))$, dominated by
$O_{\alpha,x}(|e^{\Phi_\alpha(1+i\tau)}|) \in L^1(\R)$; on $T < |\tau| \le 4k$,
Lemma~\ref{lem:W0-props}(ii),(iv) bound $|W_0((1+i\tau)/k; 1)/W_0(1/k; 1)|^k \ll L e^{-\pi T/2} \ll k^{-150}$;
on $|\tau| > 4k$, Lemma~\ref{lem:W0-props}(iii) gives $\int_{|\tau|>4k} (5k/(2|\tau|))^k \dd\tau \ll k(5/8)^k$.
Letting $k \to \infty$ gives $\nu_\alpha(x) \ge 0$ for all $x \in \R$ as the limit of the
non-negative normalized convolution density on the left-hand side.

Now set $p_\alpha(x) := e^x \nu_\alpha(x) \ge 0$ and fix $c \ge 1$. On $1 \le \re z \le c$,
$\re\Phi_\alpha(u+i\tau) \sim -\frac{\pi}{2}|\tau|$ uniformly as $|\tau|\to\infty$, so Cauchy's theorem
shifts the contour from $\re z = 1$ to $\re z = c$, giving
\[
f(x) := e^{(c-1)x} p_\alpha(x) = \frac{1}{2\pi}\int_{-\infty}^\infty g(\tau) e^{-i\tau x}\dd\tau,
\qquad g(\tau) := e^{\Phi_\alpha(c+i\tau)}.
\]
Since $\Phi_\alpha'(z) = \gamma + \alpha + 1 + \log z$ and $\Phi_\alpha''(z) = 1/z$ satisfy
$|\Phi_\alpha'(c+i\tau)| \ll_{c,\alpha} 1 + \log(1+|\tau|)$, $|\Phi_\alpha''(c+i\tau)| \le 1$, and
$|g(\tau)| \ll_{c,\alpha} e^{-\pi|\tau|/4}$, we have $g, g', g'' \in L^1(\R_\tau)$ vanishing at
$\pm\infty$. Integrating by parts twice yields $|f(x)| \le \frac{1}{2\pi}\min(\|g\|_{L^1}, x^{-2}\|g''\|_{L^1})$,
so $0 \le f(x) \le C_{c,\alpha}/(1+x^2) \in L^1(\R)$. Fourier inversion on $L^1(\R)$ at $\tau=0$
therefore gives $\int_{-\infty}^\infty p_\alpha(x)\dd x = e^{\gamma+\alpha}$ (at $c=1$) and
$\int_{-\infty}^\infty e^{(c-1)x} p_\alpha(x)\dd x = e^{c(\gamma+\alpha+\log c)} \to +\infty$ as $c \to \infty$.
If $p_\alpha$ vanished on $(0,\infty)$, the left-hand side would be $\le \int_{-\infty}^0 p_\alpha(x)\dd x \le e^{\gamma+\alpha}$
for all $c > 1$, a contradiction. By real-analyticity, the zeros of $\nu_\alpha$ in $\R$ are discrete,
so $\int_0^\infty \widetilde{Q}(x)^2 \nu_\alpha(x)\dd x > 0$.
\end{proof}

\begin{proof}[Proof of Theorem~\ref{thm:uniqueness}]
Set $\tilde{g} := c^{-1} g$, $\lambda := c^2 / m_0(g) > 0$, and $g^*(u) := \lambda \tilde{g}(\lambda u)$.

\smallskip
\noindent\textbf{Step 1: Normalization and scale invariance.}
Changing variables $v = \lambda u$ yields $m_0(g^*) = \lambda m_0(\tilde{g}) = 1$, $g^* \in L^2 \cap L^4$,
and by~(H2), $g^*(u) = u^{-1} + O_g(u^{-1-\delta})$ as $u \to \infty$. At $a = \lambda k\sigma$
and $a' = k\sigma$, we have $g(a t_i) = \frac{m_0(g)}{c} g^*(a' t_i)$, so $(m_0(g)/c)^k$ cancels
in $R_k(F_g) = k J(F_g)/I(F_g) = R_k(F_{g^*})$.

\smallskip
\noindent\textbf{Step 2: Inner-integral constant $\beta_*$ and Laplace representation.}
Because $g^*, h \in L^2_{\mathrm{loc}}[0,\infty)$ and $g^*(u) - h(u) = O_g(u^{-1-\delta})$ at $\infty$,
we have $g^* - h \in L^1(0,\infty)$ and set $\beta_* := \int_0^\infty (g^*(u) - h(u))\dd u$.
Since $(g^*)^2 \in L^1 \cap L^2$, putting $W_{g^*}(\zeta) := \int_0^\infty g^*(u)^2 e^{-\zeta u}\dd u$
($\re\zeta > 0$) gives $\omega \mapsto W_{g^*}((\sigma+i\omega)/a') \in L^2(\R)$ by Plancherel;
setting $v_{Q,g^*}(y) := \int_0^y \widetilde{Q}(\sigma(y-t)) g^*(a' t)\dd t$ and applying
Cauchy--Schwarz gives $|v_{Q,g^*}(y)| \le (\int_0^y \widetilde{Q}(\sigma(y-t))^2\dd t)^{1/2}\|g^*(a'\cdot)\|_{L^2}
\ll_{Q, a'} (1+y)^{d+1/2}$ ($d = \deg\widetilde{Q}$), ensuring absolute
integrability for $k \ge 3$ and exact Laplace representations (cf.\ \cite[Theorem~4]{Lewulis}):
\begin{align*}
I(F_{g^*}) &= \frac{1}{2\pi i}\int_{\sigma-i\infty}^{\sigma+i\infty} K_Q(s)\, W_{g^*}(s/a')^k (a')^{-k}\dd s, \\
J(F_{g^*}) &= \frac{1}{2\pi i}\int_{\sigma-i\infty}^{\sigma+i\infty} e^s W_{g^*}(s/a')^{k-1}(a')^{-(k-1)} V_{Q,g^*}(s)\dd s,
\end{align*}
where $K_Q(s) = \int_0^1 \widetilde{Q}(\sigma(1-X))^2 e^{sX}\dd X$ and
$V_{Q,g^*}(s) = \int_0^\infty \bigl(\int_0^y \widetilde{Q}(\sigma(y-t)) g^*(a' t)\dd t\bigr)^2 e^{-sy}\dd y$.
In macro coordinates $x = \sigma y$, $u = k\sigma t$, $s = \sigma z$ ($z = 1+i\tau$),
\[
\sigma^3 k^2 V_{Q, g^*}(\sigma z)
= \int_0^\infty \left(\int_0^{kx} \widetilde{Q}\!\left(x - \frac{u}{k}\right) g^*(u)\dd u\right)^{\!2} e^{-zx}\dd x.
\]
Expanding $\widetilde{Q}(x-u/k) = \widetilde{Q}(x) + \sum_{j=1}^d \frac{(-1)^j \widetilde{Q}^{(j)}(x)}{j! k^j} u^j$
and writing $g^* = h + (g^* - h)$ yields
\[
\int_0^{kx} \widetilde{Q}\!\left(x - \frac{u}{k}\right) g^*(u)\dd u
= L\,\widetilde{Q}(x) + \widetilde{Q}(x)(\log x + \beta_*) - (\cH\widetilde{Q})(x) + \widetilde{R}_k(x),
\]
where $(\cH\widetilde{Q})(x) := -\int_0^x \frac{\widetilde{Q}(x-w)-\widetilde{Q}(x)}{w}\dd w = \sum_{j=0}^d c_j \Harm_j x^j$
for $\widetilde{Q}(x) = \sum_{j=0}^d c_j x^j$. Let us verify $\|\widetilde{R}_k\|_{L^2(e^{-x}\dd x)} \to 0$:
\begin{itemize}
\item The $h$-part of $\widetilde{R}_k$ is bounded in Lemma~\ref{lem:uniform-W-V}(ii) below
(independently of Section~\ref{sec:uniqueness}) with $L^2(e^{-x}\dd x)$-norm $O_{\widetilde{Q}}(k^{-1/2})$.
\item The tail $-\widetilde{Q}(x)\int_{kx}^\infty (g^*(u)-h(u))\dd u$ tends to $0$ pointwise and is
dominated by $\|g^*-h\|_{L^1}|\widetilde{Q}(x)| \in L^2(e^{-x}\dd x)$.
\item For each $1 \le j \le d$, choose $U \ge 1$ so $|g^*(u)-h(u)| \ll_g u^{-1-\delta'}$ for $u \ge U$,
with $\delta' := \min(\delta, 1/2) \in (0, 1/2]$ (so $j - \delta' \ge 1/2$). For $k \ge U^2$ ($\sqrt{k} \ge U$),
splitting $k^{-j}\int_0^{kx} u^j |g^*-h|\dd u$ at $\min(kx, \sqrt{k})$ gives $\le k^{-1/2}\|g^*-h\|_{L^1}$ on
$[0, \min(kx, \sqrt{k})]$, while the second part on $[\min(kx, \sqrt{k}), kx]$ vanishes when $kx \le \sqrt{k}$
and is $\ll_g k^{-j}\int_{\sqrt{k}}^{kx} u^{j-1-\delta'}\dd u \le \frac{x^{j-\delta'}}{(j-\delta')k^{\delta'}}$
when $kx > \sqrt{k}$; both bounds are $o(1)$ in $L^2(e^{-x}\dd x)$.
\end{itemize}
Expanding the square against $e^{-zx}$ ($|e^{-zx}| = e^{-x}$) via Cauchy--Schwarz gives uniformly in $\tau \in \R$:
\begin{equation}\label{eq:Vg-exp}
\sigma^3 k^2 V_{Q, g^*}(\sigma z) = L^2 N_0(\tau) + 2L\bigl(N_1(\tau) + \beta_* N_0(\tau)\bigr) + o_{\widetilde{Q}, g}(L),
\end{equation}
with $N_0(\tau) := \int_0^\infty \widetilde{Q}(x)^2 e^{-zx}\dd x$ and
$N_1(\tau) := \int_0^\infty \widetilde{Q}(x)\bigl(\widetilde{Q}(x)\log x - (\cH\widetilde{Q})(x)\bigr) e^{-zx}\dd x$,
and in particular $\sigma^3 k^2 |V_{Q, g^*}(\sigma z)| \ll_{\widetilde{Q}, g} L^2$.

\smallskip
\noindent\textbf{Step 3: Uniform Laplace expansion of $W_{g^*}(\zeta)$.}
With $U \ge 1$ from Step~2 and $\delta' := \min(\delta, 1/2) \in (0, 1/2]$,
$\int_0^\infty u(g^*(u)-h(u))^2\dd u \le U\|g^*-h\|_{L^2}^2 + O_g(\int_U^\infty u^{-1-2\delta'}\dd u) < \infty$.
Since $0 \le u h(u) \le 1$ and $g^*-h \in L^1$, the identity
$u((g^*)^2 - h^2) = u(g^*-h)^2 + 2u h(u)(g^*-h)$ places $u((g^*)^2-h^2) \in L^1(0,\infty)$, so we may set
\[
c_1(g^*) := \gamma - \int_0^\infty u\bigl(g^*(u)^2 - h(u)^2\bigr)\dd u.
\]
Because $\int_0^\infty (g^*(u)^2 - h(u)^2)\dd u = m_0(g^*) - m_0(h) = 0$, for $\re\zeta > 0$ we have
\[
\frac{W_{g^*}(\zeta) - W_h(\zeta)}{\zeta} + \int_0^\infty u\bigl(g^*(u)^2 - h(u)^2\bigr)\dd u
= \int_0^\infty \bigl(g^*(u)^2 - h(u)^2\bigr)\frac{e^{-\zeta u} - 1 + \zeta u}{\zeta}\dd u.
\]
Using $|e^{-w}-1| \le |w|$ and $|e^{-w}-1+w| \le |w|^2/2$ ($\re w \ge 0$), the kernel satisfies
$|(e^{-\zeta u}-1+\zeta u)/\zeta| \le \min(2u, \frac{1}{2}|\zeta| u^2)$. Splitting at $u = |\zeta|^{-1/2}$
against the integrable majorant $2u|(g^*)^2-h^2|$ makes the integral $o(1)$ uniformly as $|\zeta| \to 0$
in $\{\re\zeta > 0\}$. Combined with $W_h(\zeta) = 1 + \zeta\log\zeta + \gamma\zeta + O(|\zeta|^2(1+|\log\zeta|))$
(Lemma~\ref{lem:W0-props} below), we obtain
\begin{equation}\label{eq:Wg-exp}
W_{g^*}(\zeta) = 1 + \zeta\log\zeta + c_1(g^*)\zeta + r_g(\zeta),
\qquad \sup_{\re\zeta > 0,\, |\zeta| \le r} \frac{|r_g(\zeta)|}{|\zeta|} \xrightarrow[r\to 0^+]{} 0.
\end{equation}

\smallskip
\noindent\textbf{Step 4: Contour passage $k \to \infty$ and shift covariance.}
Let $\zeta = (1+i\tau)/k$, $\zeta_0 = 1/k$, and split $\tau \in \R$ at $M \ge 16$ and $\delta_0 k$
with $\delta_0 \in (0, e^{-3})$ sufficiently small.

\emph{Region 1 ($|\tau| \le M$):} By~\eqref{eq:Wg-exp}, uniformly on $|\tau| \le M$,
$k\log\frac{W_{g^*}(\zeta)}{W_{g^*}(\zeta_0)} = i\tau(c_1(g^*)-L) + (1+i\tau)\log(1+i\tau) + o_M(1)$,
so multiplying by $e^{i\sigma\tau} = e^{i(L+\alpha)\tau}$ gives
\[
e^{i\sigma\tau}\left(\frac{W_{g^*}(\zeta)}{W_{g^*}(\zeta_0)}\right)^{\!k}
= e^{-(\gamma+\alpha')} e^{\Phi_{\alpha'}(1+i\tau)}\bigl(1 + o_M(1)\bigr),
\qquad \alpha' := \alpha + c_1(g^*) - \gamma.
\]

\emph{Region 2 ($M < |\tau| \le \delta_0 k$):} Choose $M \ge 16$ so $\arctan M > 5/4$ and
$\log(2t) \le \frac{1}{4} t$ ($t \ge M$). Write $\zeta = \rho e^{i\theta}$ ($\rho = \sqrt{1+\tau^2}/k \le 2\delta_0$,
$\theta = \arctan\tau$), $W_{g^*}(\zeta) = 1 + X + i Y$, and $W_{g^*}(\zeta_0) = 1 + X_0$ where
$X := \frac{\log\rho - \tau\theta + c_1(g^*)}{k} + \re r_g(\zeta)$,
$X_0 := \frac{-L + c_1(g^*)}{k} + r_g(\zeta_0)$, and
$Y := \frac{\tau\log\rho + \theta + c_1(g^*)\tau}{k} + \im r_g(\zeta)$.
Put $\varepsilon(\delta_0) := \sup_{0 < |\zeta| \le 2\delta_0,\, \re\zeta > 0} |r_g(\zeta)|/|\zeta| \to 0^+$.
Because $\log\rho + L \le \frac{1}{4}|\tau|$, $\tau\theta \ge \frac{5}{4}|\tau|$, and
$|r_g(\zeta)| + |r_g(\zeta_0)| \le 3\varepsilon(\delta_0)|\tau|/k$, taking $\delta_0$ small enough
that $\varepsilon(\delta_0) < 1/12$ gives $X - X_0 \le -\frac{3|\tau|}{4k}$, while
$|\tau\theta|/k \le \frac{\pi}{2}\delta_0 < 0.08$ (for $\delta_0 < e^{-3}$),
$|r_g(\zeta)|+|r_g(\zeta_0)| < \frac{1}{4}\delta_0$, and $\frac{|\log\rho|+L+2|c_1(g^*)|}{k} \to 0$
yield $|X| + |X_0| < 1/4$ for sufficiently large $k$, so
$(1+X)^2 - (1+X_0)^2 = (2+X+X_0)(X-X_0) \le -\frac{|\tau|}{k}$. Meanwhile, with
$u := |\tau|/k \in (0,\delta_0]$ and $C_2 := |c_1(g^*)| + \frac{\pi}{2} + 1$,
$|Y| \le \frac{2|\tau|\log(1/u) + C_2|\tau|}{k}$ implies
$Y^2 \le \frac{|\tau|}{k}\bigl(8 u\log^2(1/u) + 2 C_2^2 u\bigr) < \frac{|\tau|}{4k}$
once $\delta_0 \in (0, e^{-3})$ is chosen so $8\delta_0\log^2(1/\delta_0) + 2 C_2^2\delta_0 < 1/4$
(since $u \mapsto u\log^2(1/u)$ increases on $(0, e^{-2})$). As $0 < W_{g^*}(\zeta_0) < 1$,
\[
\left|\frac{W_{g^*}(\zeta)}{W_{g^*}(\zeta_0)}\right|^2
= \frac{(1+X)^2 + Y^2}{(1+X_0)^2} \le 1 - \frac{3|\tau|}{4k W_{g^*}(\zeta_0)^2} \le e^{-|\tau|/(2k)},
\]
whence $\bigl|W_{g^*}(\zeta)/W_{g^*}(\zeta_0)\bigr|^{k-1} \le e^{-(k-1)|\tau|/(4k)} \le e^{-|\tau|/8}$.

\emph{Region 3 ($|\tau| > \delta_0 k$):} For $\omega = \tau/k$ ($|\omega| > \delta_0$), the
Riemann--Lebesgue lemma implies that $\widehat{(g^*)^2}(\omega) = \int_0^\infty g^*(u)^2 e^{-i\omega u}\dd u$
is continuous and vanishes at $\infty$, with $|\widehat{(g^*)^2}(\omega)| < 1$ for every $\omega \neq 0$ (equality would force $e^{-i\omega u}$
to be constant a.e.\ on $\{g^* \neq 0\}$, which has positive measure while level sets of $e^{-i\omega u}$ are discrete
cosets of $\frac{2\pi}{|\omega|}\mathbb{Z}$). Thus $\theta_1 := \sup_{|\omega| \ge \delta_0}|\widehat{(g^*)^2}(\omega)| < 1$.
Since $W_{g^*}(1/k+i\omega) \to \widehat{(g^*)^2}(\omega)$ uniformly by dominated convergence and $W_{g^*}(1/k) \to 1$,
there exists $\theta_0 \in (\theta_1, 1)$ with $|W_{g^*}(\zeta)/W_{g^*}(\zeta_0)| \le \theta_0$ for large $k$
and $|\tau| > \delta_0 k$. By Plancherel, $\int_{-\infty}^\infty |W_{g^*}(1/k+i\omega)|^2\dd\omega \le 2\pi\|g^*\|_{L^4}^4$;
since the squared inner integrands are non-negative, $|K_Q(\sigma z)| \le K_Q(\sigma) \ll e^\sigma/\sigma$
and $\sigma^3 k^2 |V_{Q,g^*}(\sigma z)| \le \sigma^3 k^2 V_{Q,g^*}(\sigma) \ll L^2$; thus
Region~3 contributes at most $O_{\widetilde{Q},g}(k)\,\theta_0^{k-3}$ to the normalized Bromwich integrals
(after factoring out $e^\sigma W_{g^*}(1/k)^k (a')^{-k}$ from $I(F_{g^*})$ and
$\frac{e^\sigma L^2}{\sigma W_{g^*}(1/k)} W_{g^*}(1/k)^k (a')^{-k}$ from $k J(F_{g^*})$, and likewise in $A_k, A_k', B_k$).

To pass to the limit cleanly in the $L^2$-scaled numerator, keep the contour factor unexpanded until after
division: for $k \ge 4$, define
\[
A_k := \frac{1}{2\pi}\int_{-\infty}^\infty e^{i\sigma\tau}\!\left(\frac{W_{g^*}(\zeta)}{W_{g^*}(\zeta_0)}\right)^{\!k} N_0(\tau)\dd\tau,
\quad
A_k' := \frac{1}{2\pi}\int_{-\infty}^\infty e^{i\sigma\tau}\!\left(\frac{W_{g^*}(\zeta)}{W_{g^*}(\zeta_0)}\right)^{\!k-1} N_0(\tau)\dd\tau,
\]
and $B_k := \frac{1}{2\pi}\int_{-\infty}^\infty e^{i\sigma\tau}(W_{g^*}(\zeta)/W_{g^*}(\zeta_0))^{k-1}(N_1(\tau) + \beta_* N_0(\tau))\dd\tau$.
By~\eqref{eq:Phi-abs}, Region~2's $e^{-|\tau|/8}$, Region~3's $O(k\theta_0^{k-3})$, and Region~1 on $|\tau| \le M(\varepsilon)$
($M(\varepsilon) \to \infty$ as $\varepsilon \to 0^+$), letting $k \to \infty$ gives
\[
A_k \to e^{-(\gamma+\alpha')}\mathcal{I}_\infty(\alpha'),
\qquad
B_k \to e^{-(\gamma+\alpha')}\bigl(\mathcal{M}_\infty(\alpha') + \beta_*\mathcal{I}_\infty(\alpha')\bigr),
\]
where $\mathcal{I}_\infty(\alpha') := \int_0^\infty \widetilde{Q}(x)^2 \nu_{\alpha'}(x)\dd x > 0$ (Lemma~\ref{lem:nu-pos})
and $\mathcal{M}_\infty(\alpha') := \frac{1}{2\pi}\int_{-\infty}^\infty e^{\Phi_{\alpha'}(1+i\tau)} N_1(\tau)\dd\tau$.
In particular, $A_k$ is bounded away from $0$ for large $k$. Moreover, on $|\tau| \le \delta_0 k$, \eqref{eq:Wg-exp}
yields $|1 - W_{g^*}(\zeta)/W_{g^*}(\zeta_0)| \ll_g \frac{L(1+|\tau|)}{k}$; integrating against $|W_{g^*}(\zeta)/W_{g^*}(\zeta_0)|^{k-1}|N_0(\tau)|$
across $|\tau| \le 16$, $16 < |\tau| \le \delta_0 k$ (taking $M = 16$ in Region~2), and $|\tau| > \delta_0 k$ yields $|A_k' - A_k| \ll_{\widetilde{Q},g} L/k$.

Inserting $\sigma e^{-\sigma z} K_Q(\sigma z) = N_0(\tau) + O(L^{2d}/k)$ (Lemma~\ref{lem:uniform-W-V}(iii)) and~\eqref{eq:Vg-exp}
into the Bromwich integrals (using $\int_{-\infty}^\infty |W_{g^*}(\zeta)/W_{g^*}(\zeta_0)|^{k-1}\dd\tau \ll_g 1$) gives
\begin{align*}
I(F_{g^*}) &= e^\sigma W_{g^*}(1/k)^k (a')^{-k} \bigl(A_k + O(L^{2d}/k)\bigr), \\
k J(F_{g^*}) &= \frac{e^\sigma}{\sigma W_{g^*}(1/k)} W_{g^*}(1/k)^k (a')^{-k} \bigl(L^2 A_k + 2L B_k + o(L)\bigr),
\end{align*}
since $L^2 A_k' = L^2 A_k + O(L^3/k)$. Using $(A_k + O(L^{2d}/k))^{-1} = A_k^{-1}(1 + O(L^{2d}/k))$ and
$W_{g^*}(1/k) = 1 + O_g(L/k)$, dividing yields
\[
R_k(F_{g^*}) = \frac{L^2}{\sigma} + \frac{2 B_k}{A_k} + o(1)
= \log k + \bigl(2\beta_* + c_1(g^*) - \gamma\bigr) + \eta(h; \alpha', \widetilde{Q}) + o(1),
\]
where $\eta(h; \alpha', \widetilde{Q}) = -\alpha' + 2\mathcal{M}_\infty(\alpha')/\mathcal{I}_\infty(\alpha')$.
Taking the supremum over $\alpha \in \R$ yields~\eqref{eq:shift-covariance} with $\kappa(g^*) - \gamma := 2\beta_* + c_1(g^*) - \gamma$.

\smallskip
\noindent\textbf{Step 5: Algebraic completion.}
Using $(1+u)h(u) = 1$ ($u \ge 0$) and $\int_0^\infty (g^*(u)^2 - h(u)^2)\dd u = 0$,
\begin{align*}
-\int_0^\infty (1+u)\bigl(g^*(u) - h(u)\bigr)^2 \dd u
&= -\int_0^\infty u\bigl(g^*(u)^2 - h(u)^2\bigr)\dd u + 2\int_0^\infty \bigl(g^*(u) - h(u)\bigr)\dd u \\
&= \bigl(c_1(g^*) - \gamma\bigr) + 2\beta_* = \kappa(g^*) - \gamma \le 0.
\end{align*}
Since $1+u \ge 1 > 0$, equality holds if and only if $g^*(u) = h(u)$ a.e.\ on $(0,\infty)$,
equivalently $g(u) = C/(1+\theta u)$ a.e.\ with $C = m_0(g)/c$ and $\theta = m_0(g)/c^2 > 0$.
\end{proof}

\begin{remark}[Scale choice, stability, and scope]
Rescaling $a' = \mu k\sigma$ ($\mu > 0$) in~\eqref{eq:Wg-exp} yields
$k\log\bigl(W_{g^*}(\sigma(1+i\tau)/a')/W_{g^*}(\sigma/a')\bigr) = -\frac{i\tau}{\mu}\log k + O(1)$,
which cancels $e^{i\sigma\tau}$ precisely when $\mu = 1$. Quantitative stability also follows directly
from~\eqref{eq:shift-covariance}--\eqref{eq:kappa-penalty}: any kernel $g$ coming within $\varepsilon > 0$
of the optimal radial shifted limit satisfies $\|g^* - h\|_{L^2((1+u)\dd u)} \le \sqrt{\varepsilon}$.
Theorem~\ref{thm:uniqueness} applies strictly to product kernels~\eqref{eq:Fg-def} with radial polynomials;
Theorem~\ref{thm:main} escapes this radial barrier via the non-radial power-sum ansatz~\eqref{eq:ansatz-C}.
\end{remark}

\section{Exact Finite-\texorpdfstring{$k$}{k} Laplace Reduction for the \texorpdfstring{$\Lambda_W$}{Lambda\_W} Family}\label{sec:laplace}

Let $\Lambda_W$ ($W \ge 0$) be the set of all partitions $\lambda = (\lambda_1,\dots,\lambda_{\ell(\lambda)})$
with $\lambda_1 \ge \cdots \ge \lambda_{\ell(\lambda)} \ge 2$ and $|\lambda| \le W$ (including $()$); note
$\ell(\lambda) \le \lfloor W/2 \rfloor$ ($\le 4$ for $W=9$). For parameters $a, \sigma > 0$ and polynomials
$Q_\lambda(y) := \widetilde{Q}_\lambda(\sigma y)$, set $g(t) := (1+at)^{-1}$, $w(t) := g(t)^2$, and
\begin{equation}\label{eq:F-LambdaW}
F(t_1,\dots,t_k) := \left[\sum_{\lambda \in \Lambda_W} Q_\lambda(1 - P_1)
\prod_{j=1}^{\ell(\lambda)} \bigl(\sigma^{\lambda_j} P_{\lambda_j}\bigr)\right]
\prod_{i=1}^k g(t_i).
\end{equation}
For $\re s > 0$ and integers $m \ge 0$, define
\begin{equation}\label{eq:Wm-E-def}
W_m(s; a) := \int_0^\infty \frac{t^m e^{-st}}{(1+at)^2}\dd t,
\qquad
E(s; a) := \int_0^\infty \frac{e^{-st}}{1+at}\dd t.
\end{equation}

\begin{lemma}[Partial-Fraction Representation and Decay of $W_m$]\label{lem:Wm-pf}
Let $a > 0$ and $\re s > 0$.
\begin{enumerate}
\item[(i)] $W_0(s;a) = \frac{1}{a} - \frac{s}{a} E(s;a)$, and for every integer $m \ge 2$,
\begin{equation}\label{eq:Wm-pf}
W_m(s;a) = \frac{1}{a^m}\left(\sum_{r=0}^{m-2} (-1)^{m-r}(m-1-r)\frac{r!\,a^r}{s^{r+1}}
- (-1)^m m\, E(s;a) + (-1)^m W_0(s;a)\right).
\end{equation}
\item[(ii)] For every $m \ge 0$ and fixed $\sigma > 0$, $|W_m(\sigma + i\omega; a)| \ll_{m,a,\sigma} (1 + |\omega|)^{-1}$
uniformly in $\omega \in \R$.
\end{enumerate}
\end{lemma}

\begin{proof}
Integration by parts gives $W_0(s;a) = \frac{1}{a} - \frac{s}{a} E(s;a)$. For $m \ge 2$, substituting $x = a t$
into the exact polynomial identity
\[
x^m = (1+x)^2 \sum_{r=0}^{m-2} (-1)^{m-r}(m-1-r) x^r - (-1)^m m(1+x) + (-1)^m
\]
(verified directly by comparing coefficients of $x^j$ for $0 \le j \le m$) and dividing by $(1+x)^2$ against
$a^{-m-1} e^{-(s/a)x}\dd x$ yields~\eqref{eq:Wm-pf}, using $\int_0^\infty x^r e^{-(s/a)x}\dd x = r!(a/s)^{r+1}$.
Part~(ii) follows by integrating~\eqref{eq:Wm-E-def} by parts once against $e^{-i\omega t}$, since
$t \mapsto t^m (1+at)^{-2} e^{-\sigma t} \in C^1[0,\infty) \cap L^1(0,\infty)$ has integrable derivative
vanishing at $\infty$.
\end{proof}

For any tuple $\xi = (\xi_1, \dots, \xi_n)$ of parts $\ge 2$, let $\Pi([n])$ be the set of set partitions
$\pi$ of $[n] := \{1,\dots,n\}$ ($\Pi(\emptyset) = \{\emptyset\}$), put $s_\xi(B) := \sum_{j \in B} \xi_j$
for each block $B \in \pi$, let $(k)_t := k(k-1)\cdots(k-t+1)$ ($(k)_0 := 1$), and define
\begin{equation}\label{eq:Gk-def}
\cG_k(s; \xi) := \sum_{\substack{\pi \in \Pi([\ell(\xi)]) \\ |\pi| \le k}} (k)_{|\pi|}\,
W_0(s; a)^{k - |\pi|} \prod_{B \in \pi} W_{s_\xi(B)}(s; a)
\end{equation}
(where summands with $|\pi| > k$ vanish identically due to the falling factorial $(k)_{|\pi|} = 0$,
avoiding formal negative powers of $W_0$ when $k$ is small).
For $\lambda \in \Lambda_W$ and $S \subseteq [\ell(\lambda)]$, write $s_\lambda(S) := \sum_{j \in S} \lambda_j$,
and let $\lambda_{\setminus S} \in \Lambda_W$ be the partition obtained by removing indices in $S$.

\begin{proposition}[Exact Finite-$k$ Set-Partition Laplace Formulas]\label{prop:exact-laplace}
Fix $W \ge 0$, $k \ge 3$, $a, \sigma > 0$, and polynomials $\{Q_\lambda\}_{\lambda \in \Lambda_W}$.
Define $K_{\lambda_1,\lambda_2}(s) := \int_0^1 Q_{\lambda_1}(1-X) Q_{\lambda_2}(1-X) e^{sX}\dd X$ and
\begin{equation}\label{eq:v-eff-def}
v_\mu^{\mathrm{eff}}(y) := \int_0^y \frac{Q_\mu(y-t)}{1+at}\dd t
+ \sum_{\lambda \in \Lambda_W}\sum_{\substack{\emptyset \neq S \subseteq [\ell(\lambda)] \\ \lambda_{\setminus S} = \mu}}
\sigma^{s_\lambda(S)} \int_0^y \frac{Q_\lambda(y-t)\, t^{s_\lambda(S)}}{1+at}\dd t,
\end{equation}
with $V_{\mu_1,\mu_2}^{\mathrm{eff}}(s) := \int_0^\infty v_{\mu_1}^{\mathrm{eff}}(y) v_{\mu_2}^{\mathrm{eff}}(y) e^{-sy}\dd y$.
Then
\begin{align}
I(F) &= \frac{1}{2\pi i}\int_{\sigma-i\infty}^{\sigma+i\infty}
\sum_{\lambda_1,\lambda_2 \in \Lambda_W} K_{\lambda_1,\lambda_2}(s)\,
\sigma^{|\lambda_1|+|\lambda_2|} \, \cG_k(s; \lambda_1 \sqcup \lambda_2)\dd s,
\label{eq:I-exact-laplace}\\
J(F) &= \frac{1}{2\pi i}\int_{\sigma-i\infty}^{\sigma+i\infty}
e^s \sum_{\mu_1,\mu_2 \in \Lambda_W} V_{\mu_1,\mu_2}^{\mathrm{eff}}(s)\,
\sigma^{|\mu_1|+|\mu_2|} \, \cG_{k-1}(s; \mu_1 \sqcup \mu_2)\dd s,
\label{eq:J-exact-laplace}
\end{align}
where $\lambda_1 \sqcup \lambda_2$ denotes concatenation and both vertical integrals converge absolutely.
\end{proposition}

\begin{proof}
Expanding $\prod_{j=1}^n P_{\xi_j}(t) = \sum_{i_1,\dots,i_n \in [k]} \prod_{j=1}^n t_{i_j}^{\xi_j}$
according to the equality partition $\pi \in \Pi([n])$ of indices $(i_1,\dots,i_n)$, there are
$(k)_{|\pi|}$ injective assignments of the $|\pi|$ blocks to $[k]$ (vanishing when $|\pi| > k$).
Multiplying by $\prod_{i=1}^k w(t_i) e^{-s t_i}$ and integrating over $[0,\infty)^k$ gives $\cG_k(s; \xi)$
as the Laplace transform of the hyperplane integral
\[
\Psi_k(X; \xi) := \int_{\substack{t_1,\dots,t_k \ge 0 \\ t_1+\cdots+t_k = X}}
\left(\prod_{j=1}^n P_{\xi_j}(t)\right) \prod_{i=1}^k w(t_i)\dd m_{k-1},
\]
where $\dd m_{k-1} := \dd t_1\cdots\dd t_{k-1}$ denotes the projected $(k-1)$-dimensional Lebesgue
measure normalized by $\dd X \dd m_{k-1} = \dd t_1\cdots\dd t_k$ (with $t_k = X - \sum_{i=1}^{k-1} t_i$).
By Lemma~\ref{lem:Wm-pf}(ii), each non-vanishing term ($|\pi| \le k$) in~\eqref{eq:Gk-def} is
a product of $k \ge 3$ factors $W_m$, hence $O((1+|\omega|)^{-k}) \in L^1(\R_\omega)$. Since $\Psi_k(X;\xi) = 0$
for $X < 0$ and the inversion integral is continuous on $\R$ (so $\Psi_k(0;\xi) = 0$), Fourier/Laplace inversion
gives $\Psi_k(X;\xi) = \frac{1}{2\pi i}\int_{\sigma-i\infty}^{\sigma+i\infty} e^{sX}\cG_k(s;\xi)\dd s$ for all $X \in \R$.
Integrating against $Q_{\lambda_1}(1-X) Q_{\lambda_2}(1-X)$ over $X \in [0,1]$ and applying Fubini yields~\eqref{eq:I-exact-laplace}.

For $J(F)$, setting $X := \sum_{i=1}^{k-1} t_i$, $y := 1 - X$, and $P_r' := \sum_{i=1}^{k-1} t_i^r$
($P_r = P_r' + t_k^r$), binomial expansion of $\prod_{j=1}^{\ell(\lambda)} (P_{\lambda_j}' + t_k^{\lambda_j})$
over subsets $S \subseteq [\ell(\lambda)]$ collects the inner $t_k$-integral as
$v(t_1,\dots,t_{k-1}) = \sum_{\mu \in \Lambda_W} \bigl(\prod_{r=1}^{\ell(\mu)} \sigma^{\mu_r} P_{\mu_r}'\bigr) v_\mu^{\mathrm{eff}}(y)$.
Squaring and integrating over $t_1 + \cdots + t_{k-1} = X \in [0,1]$ gives
$J(F) = \int_0^1 \sum_{\mu_1,\mu_2} v_{\mu_1}^{\mathrm{eff}}(1-X) v_{\mu_2}^{\mathrm{eff}}(1-X)
\sigma^{|\mu_1|+|\mu_2|} \Psi_{k-1}(X; \mu_1\sqcup\mu_2)\dd X$.
Since $v_\mu^{\mathrm{eff}}(y)$ has polynomial growth and vanishes at $0$, $V_{\mu_1,\mu_2}^{\mathrm{eff}}(\sigma+i\omega) \ll (1+|\omega|)^{-1}$,
while $|\cG_{k-1}(\sigma+i\omega;\xi)| \ll (1+|\omega|)^{-(k-1)} \le (1+|\omega|)^{-2}$, which belongs to
$L^1(\R_\omega)$ since $k - 1 \ge 2$.
Because $\Psi_{k-1}(1-y;\xi) = 0$ for $y > 1$, substituting $X = 1-y$ extends $\int_0^1$ to $\int_0^\infty$, and Fubini's theorem yields~\eqref{eq:J-exact-laplace}.
\end{proof}

\section{Proof of the Asymptotic Limit (\texorpdfstring{$k \to \infty$}{k -> infinity})}\label{sec:asymptotic}

Specialize Proposition~\ref{prop:exact-laplace} to $\alpha \in \R$, $L = \log k$, $\sigma = L+\alpha$,
$a = k\sigma$ (for $k \ge \max(20, 2\lfloor W/2\rfloor + 2)$ with $\sigma > 0$), fixed $\{\widetilde{Q}_\lambda\}_{\lambda\in\Lambda_W}$
of degrees $d_\lambda$, and contour $s = \sigma z$ ($z = 1+i\tau$, $\zeta = s/a = (1+i\tau)/k$, $\zeta_0 = 1/k$).

\begin{lemma}[Self-Contained Properties of $W_0$]\label{lem:W0-props}
Fix $a, \sigma > 0$.
\begin{enumerate}
\item[(i)] $W_0(z; a) = a^{-1} W_0(z/a; 1)$ for $\re z > 0$, and $W_0(\zeta; 1) = 1 - \zeta e^\zeta E_1(\zeta)$
where $E_1(\zeta) := \int_\zeta^\infty t^{-1} e^{-t}\dd t$.
\item[(ii)] $\omega \mapsto |W_0(\sigma + i\omega; a)|$ is even and strictly decreasing on $[0,\infty)$ (cf.\ \cite[Lemma~5]{Lewulis}).
\item[(iii)] Whenever $0 < \zeta_0 := \re\zeta \le 1/20$, $|W_0(\zeta;1)/W_0(\zeta_0;1)| \le \frac{5}{2|\zeta|}$ (cf.\ \cite[Lemma~6]{Lewulis}).
\item[(iv)] For $|\tau| \le T := 100 L$ and sufficiently large $k$, $W_0(1/k;1) = 1 + O(L/k)$ and
\begin{equation}\label{eq:W0-log-exp}
k \log\frac{W_0(\zeta;1)}{W_0(\zeta_0;1)} = i\tau(\gamma - L) + (1+i\tau)\log(1+i\tau) + \Delta_k(\tau),
\qquad |\Delta_k(\tau)| \ll \frac{L^2(1+\tau^2)}{k}.
\end{equation}
\end{enumerate}
\end{lemma}

\begin{proof}
Part~(i) follows by $u = a t$, integration by parts, and $w = \zeta(1+t)$.
For~(ii), evenness holds by conjugation. Writing $s = \sigma + i\omega$, Lemma~\ref{lem:Wm-pf}(i) gives
$W_0(s;a) = \frac{1}{a} - \frac{s}{a}E(s;a)$; rotating $u = st$ onto $[0,\infty)$ in $E(s;a)$ yields
$W_0(s;a) = \frac{s}{a}\int_0^\infty (\frac{1}{s} - \frac{1}{s+au})e^{-u}\dd u = \int_0^\infty \frac{u e^{-u}}{r_u + i\omega}\dd u$
($r_u := \sigma + a u > 0$; in particular $\re W_0 > 0$).
Symmetrizing over $(u,v) \in (0,\infty)^2$ cancels $\pm i\omega(r_u - r_v)$:
\[
|W_0(\sigma+i\omega; a)|^2
= \int_0^\infty\!\!\int_0^\infty \frac{u v e^{-u-v}}{r_u + r_v}
\left(\frac{r_u}{r_u^2+\omega^2} + \frac{r_v}{r_v^2+\omega^2}\right)\dd u\dd v,
\]
whose $\omega$-derivative is strictly negative for $\omega > 0$.

For~(iii), opposite integration by parts yields $W_0(\zeta;1) = \frac{1}{\zeta} - \frac{2}{\zeta}\int_0^\infty \frac{e^{-\zeta t}}{(1+t)^3}\dd t$,
so $|W_0(\zeta;1)| \le \frac{2}{|\zeta|}$. On $\zeta_0 \in (0, 1/20]$,
$E_1(\zeta_0) \le \log(\zeta_0^{-1}) + e^{-1}$, and $\zeta_0 \mapsto \zeta_0 e^{\zeta_0}(\log(\zeta_0^{-1})+e^{-1})$
increases on $(0, 1/20]$ (derivative $e^{\zeta_0}((1+\zeta_0)(\log(\zeta_0^{-1})+e^{-1}) - 1) > 0$), giving
$W_0(\zeta_0;1) \ge 1 - \frac{1}{20}e^{1/20}(\log 20 + e^{-1}) > 4/5$, hence the claim.

For~(iv), $E_1(\zeta) = -\gamma - \log\zeta - \sum_{n\ge 1} \frac{(-\zeta)^n}{n\cdot n!}$ implies
$\log W_0(\zeta;1) = \zeta(\gamma + \log\zeta) + O(|\zeta|^2(1+|\log\zeta|)^2)$ for $|\zeta| \le c_0$
(with $c_0 > 0$ a sufficiently small absolute constant ensuring $|W_0(\zeta;1) - 1| \le 1/2$).
Evaluating at $\zeta = (1+i\tau)/k$ ($|\tau| \le T$) and subtracting $\zeta_0 = 1/k$ yields~\eqref{eq:W0-log-exp}.
\end{proof}

\begin{lemma}[Uniform Moment and Inner-Integral Estimates]\label{lem:uniform-W-V}
Let $\sigma = L + \alpha$, $a = k\sigma$, and $z = 1+i\tau$ ($\tau \in \R$).
\begin{enumerate}
\item[(i)] For every integer $m \ge 2$ and uniformly over all $\tau \in \R$,
\begin{equation}\label{eq:Wm-uniform}
k \sigma^m \frac{W_m(\sigma z; k\sigma)}{W_0(\sigma; k\sigma)} = \frac{(m-2)!}{z^{m-1}} + O_m\!\left(\frac{L}{k}\right).
\end{equation}
\item[(ii)] Define $(\cH\widetilde{Q})(x) := -\int_0^x \frac{\widetilde{Q}(x-u)-\widetilde{Q}(x)}{u}\dd u$,
$(B_{s-1}\widetilde{Q})(x) := \int_0^x \widetilde{Q}(x-u) u^{s-1}\dd u$ ($s \ge 2$), and
\[
G_\mu(x) := \widetilde{Q}_\mu(x)\log x - (\cH\widetilde{Q}_\mu)(x)
+ \sum_{\lambda\in\Lambda_W}\sum_{\substack{\emptyset\neq S\subseteq[\ell(\lambda)] \\ \lambda_{\setminus S}=\mu}}
\bigl(B_{s_\lambda(S)-1}\widetilde{Q}_\lambda\bigr)(x).
\]
Then uniformly over all $\tau \in \R$,
\begin{equation}\label{eq:V-uniform}
\sigma^3 k^2 V_{\mu_1,\mu_2}^{\mathrm{eff}}(\sigma z; k\sigma)
= L^2 N_{\mu_1,\mu_2}^{(0)}(\tau)
+ L\bigl(N_{\mu_1,\mu_2}^{(1)}(\tau) + N_{\mu_2,\mu_1}^{(1)}(\tau)\bigr)
+ N_{\mu_1,\mu_2}^{(2)}(\tau) + O\!\left(\frac{L^2}{k}\right),
\end{equation}
where $N_{\mu_1,\mu_2}^{(0)}(\tau) := \int_0^\infty \widetilde{Q}_{\mu_1}\widetilde{Q}_{\mu_2} e^{-zx}\dd x$,
$N_{\mu_1,\mu_2}^{(1)}(\tau) := \int_0^\infty \widetilde{Q}_{\mu_1} G_{\mu_2} e^{-zx}\dd x$, and
$N_{\mu_1,\mu_2}^{(2)}(\tau) := \int_0^\infty G_{\mu_1} G_{\mu_2} e^{-zx}\dd x$ are uniformly $O(1)$ in $\tau \in \R$.
\item[(iii)] Uniformly over $\tau \in \R$, $\sigma e^{-\sigma z} K_{\lambda_1,\lambda_2}(\sigma z)
= N_{\lambda_1,\lambda_2}^{(0)}(\tau) + O(L^{d_{\lambda_1}+d_{\lambda_2}}/k)$.
\end{enumerate}
\end{lemma}

\begin{proof}
For~(i), substituting $w = \sigma t$ expresses the ratio as
$\frac{1}{W_0(1/k;1)}\int_0^\infty \frac{k^2 w^m}{(1+kw)^2} e^{-zw}\dd w$.
Since $0 \le w^{m-2} - \frac{k^2 w^m}{(1+kw)^2} = w^{m-2}\frac{2kw+1}{(1+kw)^2}$ and $|e^{-zw}| = e^{-w}$,
splitting at $w=1$ when $m=2$ bounds the difference against $\int_0^\infty w^{m-2}e^{-zw}\dd w = \frac{(m-2)!}{z^{m-1}}$
by $O(L/k)$, while for $m \ge 3$ using $w\frac{2kw+1}{(1+kw)^2} \le \frac{2}{k}$ yields $O_m(1/k)$. Dividing
by $W_0(1/k;1) = 1 + O(L/k)$ gives~\eqref{eq:Wm-uniform}.

For~(ii), putting $x = \sigma y$ and $u = \sigma t$ in~\eqref{eq:v-eff-def} and expanding $\widetilde{Q}_\mu(x-u)$
via $\int_0^x \frac{\dd u}{u+1/k} = L + \log x + \log(1 + \frac{1}{kx})$ and
$\int_0^x \frac{u^j\dd u}{u+1/k} = \frac{x^j}{j} - \frac{1}{k}\int_0^x \frac{u^{j-1}\dd u}{u+1/k}$ gives
$\sigma k v_\mu^{\mathrm{eff}}(x/\sigma) = L \widetilde{Q}_\mu(x) + G_\mu(x) + E_{k,\mu}(x)$, where
$E_{k,\mu}$ combines $R_{1,k}^{\mu,(0)}(x) := \widetilde{Q}_\mu(x)\log(1 + \frac{1}{kx})$,
$R_{1,k}^{\mu,(1)}(x) := -\sum_{j=1}^{d_\mu} \frac{(-1)^j \widetilde{Q}_\mu^{(j)}(x)}{j!\, k}\int_0^x \frac{u^{j-1}\dd u}{u+1/k}$,
and the shift remainders $R_{2,k}^{\lambda,S}$ bounded in modulus by $\frac{1}{k}\int_0^x |\widetilde{Q}_\lambda(x-u)|u^{s-2}\dd u \ll (1+x^{d_\lambda+s-1})/k$.
Set $d_{\max} := \max_\lambda (d_\lambda + |\lambda|)$. For $R_{1,k}^{\mu,(1)}$, since
$\int_0^x \frac{u^{j-1}\dd u}{u+1/k}$ equals $\log(kx+1) \ll L + \log(1+x)$ for $j=1$ and at most
$x^{j-1}/(j-1)$ for $j \ge 2$, we have $|R_{1,k}^{\mu,(1)}(x)| \ll (1+x^{d_{\max}})\frac{L+\log(1+x)}{k}$,
contributing $O(L/k)$ in $L^1$ and $L^2$ and $O(L^2/k)$ against $A_{k,\mu}$. Splitting $(0,\infty)$
at $1/k$ and $1$ for $R_{1,k}^{\mu,(0)}$, substituting $v = kx$ on $(0, 1/k]$ gives
$\int_0^{1/k} (1+x^{d_{\max}})\log^r(1+\frac{1}{kx})e^{-x}\dd x \ll 1/k$ ($r = 1,2$),
while $\log(1+\frac{1}{kx}) \le \frac{1}{kx}$ on $(1/k,\infty)$ yields
$\int_{1/k}^\infty (1+x^{d_{\max}})\log^r(1+\frac{1}{kx})e^{-x}\dd x \ll L/k$ ($r = 1$) and $O(1/k)$ ($r = 2$),
whence $\|E_{k,\mu}\|_{L^1(e^{-x}\dd x)} \ll L/k$ and $\|E_{k,\mu}\|_{L^2(e^{-x}\dd x)} \ll k^{-1/2}$.
For cross terms with $A_{k,\mu} := L\widetilde{Q}_\mu + G_\mu$ (where $|A_{k,\mu}(x)| \ll (1+x^{d_{\max}})(L + |\log x|)$),
integrating over $(0, 1/k] \cup (1/k, 1] \cup (1,\infty)$ bounds every piece by $O(L^2/k)$, proving~\eqref{eq:V-uniform}.

For~(iii), substituting $x = \sigma(1-X)$ turns $\sigma e^{-\sigma z} K_{\lambda_1,\lambda_2}(\sigma z)$ into
$N_{\lambda_1,\lambda_2}^{(0)}(\tau) - \int_\sigma^\infty \widetilde{Q}_{\lambda_1}\widetilde{Q}_{\lambda_2} e^{-zx}\dd x$,
with tail $\ll \sigma^{d_{\lambda_1}+d_{\lambda_2}} e^{-\sigma} \ll L^{d_{\lambda_1}+d_{\lambda_2}}/k$.
\end{proof}

For any tuple $\xi = (\xi_1,\dots,\xi_n)$ of parts $\ge 2$, define the Poisson-point-process (PPP) weight
\begin{equation}\label{eq:WPPP-def}
W_\xi^{\mathrm{PPP}}(z) := \sum_{\pi \in \Pi([n])} \prod_{B \in \pi}
\frac{(s_\xi(B)-2)!}{z^{s_\xi(B)-1}} =: \sum_{r \ge 0} w_{\xi, r} z^{-r},
\end{equation}
noting $w_{(),0} = 1$, $w_{\xi, r} > 0$ iff $|\xi| - n \le r \le |\xi| - 1$ ($n \ge 1$), and
$|W_\xi^{\mathrm{PPP}}(1+i\tau)| \le W_\xi^{\mathrm{PPP}}(1)$.

\begin{proposition}[Asymptotic Contour Limit]\label{prop:contour-limit}
Let $F$ be given by~\eqref{eq:F-LambdaW} at $a = k\sigma$ ($\sigma = L+\alpha$). Define
\begin{align}
\mathcal{I}_\infty(\alpha) &:= \frac{1}{2\pi}\int_{-\infty}^\infty e^{\Phi_\alpha(1+i\tau)}
\sum_{\lambda_1,\lambda_2 \in \Lambda_W} W_{\lambda_1\sqcup\lambda_2}^{\mathrm{PPP}}(1+i\tau)\, N_{\lambda_1,\lambda_2}^{(0)}(\tau)\dd\tau, \label{eq:I-inf-def}\\
\mathcal{M}_\infty(\alpha) &:= \frac{1}{2\pi}\int_{-\infty}^\infty e^{\Phi_\alpha(1+i\tau)}
\sum_{\mu_1,\mu_2 \in \Lambda_W} W_{\mu_1\sqcup\mu_2}^{\mathrm{PPP}}(1+i\tau)\, N_{\mu_1,\mu_2}^{(1)}(\tau)\dd\tau. \label{eq:M-inf-def}
\end{align}
If $\mathcal{I}_\infty(\alpha) > 0$, then as $k \to \infty$,
\begin{equation}\label{eq:quotient-limit}
R_k(F) = L + \left(\frac{2\mathcal{M}_\infty(\alpha)}{\mathcal{I}_\infty(\alpha)} - \alpha\right) + O\!\left(\frac{1}{L}\right).
\end{equation}
\end{proposition}

\begin{proof}
Write $I(F) = e^\sigma W_0(\sigma; a)^k Z_{I,k}$ and $J(F) = \frac{e^\sigma W_0(\sigma; a)^{k-1}}{\sigma^2 k^2} Z_{J,k}$
in~\eqref{eq:I-exact-laplace}--\eqref{eq:J-exact-laplace}, with $s = \sigma z$ ($z = 1+i\tau$) and
$\widehat{\cG}_n(\sigma z;\xi) := \cG_n(\sigma z;\xi)/W_0(\sigma;a)^n$ ($n \in \{k, k-1\}$). Split $\tau \in \R$
at $T := 100 L$ and $4k$, and set $p_{\max} := 2\lfloor W/2\rfloor$ ($\le 8$ for $W=9$).

On $|\tau| > T$, $|W_m(\sigma z; a)| \le W_m(\sigma; a)$, $|K_{\lambda_1,\lambda_2}(\sigma z)| \ll e^\sigma/\sigma$,
$|V_{\mu_1,\mu_2}^{\mathrm{eff}}(\sigma z; a)| \ll L^2/(\sigma^3 k^2)$, and by Lemma~\ref{lem:uniform-W-V}(i)
each block satisfies $k\sigma^{s_\xi(B)}|W_{s_\xi(B)}|/W_0(\sigma;a) \ll 1$. Because $(n)_{|\pi|}/k^{|\pi|} \le 1$
and $n - |\pi| \ge k - 1 - p_{\max}$ ($\ge k-9$ for $W=9$), the integrands are bounded by
$O(L^2)|W_0(\sigma z;a)/W_0(\sigma;a)|^{k-1-p_{\max}}$. By Lemma~\ref{lem:W0-props}(ii),(iv), on $T < |\tau| \le 4k$
this is $\ll L^3 e^{-\pi T/2} \ll k^{-150}$ (integrating to $O(k^{-148})$), and on $|\tau| > 4k$ Lemma~\ref{lem:W0-props}(iii)
gives $\ll L^2 k (5/8)^{k-1-p_{\max}}$.

On $|\tau| \le T$, $W_0(\sigma z;a) \neq 0$, $(n)_{|\pi|}/k^{|\pi|} = 1 + O(1/k)$, and
Lemma~\ref{lem:W0-props}(iv) plus~\eqref{eq:Wm-uniform} yield, for $n \in \{k, k-1\}$,
\begin{align*}
\sigma^{|\xi|}\widehat{\cG}_n(\sigma z;\xi)
&= \left(\frac{W_0(\sigma z;a)}{W_0(\sigma;a)}\right)^{\!n}
\left(W_\xi^{\mathrm{PPP}}(1+i\tau) + O\!\left(\frac{L(1+|\tau|)}{k}\right)\right), \\
e^{i\sigma\tau}\left(\frac{W_0(\sigma z;a)}{W_0(\sigma;a)}\right)^{\!n}
&= e^{-(\gamma+\alpha)} e^{\Phi_\alpha(1+i\tau)}
+ O\!\left(\bigl|e^{\Phi_\alpha(1+i\tau)}\bigr|\frac{L^2(1+\tau^2)}{k}\right).
\end{align*}
Combining these expansions with Lemma~\ref{lem:uniform-W-V}(ii)--(iii) and~\eqref{eq:Phi-abs} yields
\begin{align*}
Z_{I,k} &= e^{-(\gamma+\alpha)}\bigl(\mathcal{I}_\infty(\alpha) + O(L^r/k)\bigr), \\
Z_{J,k} &= e^{-(\gamma+\alpha)}\bigl(L^2\mathcal{I}_\infty(\alpha) + 2L\mathcal{M}_\infty(\alpha) + O(1)\bigr),
\end{align*}
where $r := \max(2, \max_\lambda 2d_\lambda)$ and symmetry in $(\mu_1,\mu_2)$ collects the
linear-$L$ terms into $2\mathcal{M}_\infty(\alpha)$. Dividing via $k W_0(\sigma; a) = \sigma^{-1}(1 + O(L/k))$
establishes~\eqref{eq:quotient-limit}.
\end{proof}

Now define the one-dimensional moments (for integers $n \ge 0$ and $\alpha \in \R$)
\begin{align}
U_n(\alpha) &:= \frac{n!}{\pi}\re \int_0^\infty \frac{e^{\Phi_\alpha(1+i\tau)}}{(1+i\tau)^{n+1}}\dd\tau, \nonumber\\
L_n(\alpha) &:= \frac{n!}{\pi}\re \int_0^\infty \bigl(\gamma + \log(1+i\tau)\bigr)\frac{e^{\Phi_\alpha(1+i\tau)}}{(1+i\tau)^{n+1}}\dd\tau. \label{eq:Un-Ln-def}
\end{align}

\begin{proposition}[Moment Reduction and Integration-by-Parts Identity]\label{prop:UL-reduction}
For all integers $n \ge 0$ and $\alpha \in \R$:
\begin{enumerate}
\item[(i)] $L_n(\alpha) = U_{n+1}(\alpha) - (1+\alpha) U_n(\alpha)$.
\item[(ii)] Writing $\widetilde{Q}_\lambda(x) = \sum_{i=0}^{d_\lambda} c_{\lambda, i} x^i$ with
stacked vector $c \in \R^N$ ($N = \sum_{\lambda\in\Lambda_W}(d_\lambda + 1)$), we have
$\mathcal{I}_\infty(\alpha) = c^\top I(\alpha) c$ and
$2\mathcal{M}_\infty(\alpha) - \alpha \mathcal{I}_\infty(\alpha) = c^\top D_{\mathrm{net}}(\alpha) c$,
where $D_{\mathrm{net}}(\alpha) := M(\alpha) + M(\alpha)^\top - \alpha I(\alpha)$ and, for basis
indices $(\lambda_1, i), (\lambda_2, j)$ with $n := i+j$,
\begin{align}
I_{(\lambda_1,i),(\lambda_2,j)}(\alpha)
&= \sum_r w_{\lambda_1\sqcup\lambda_2,\, r} \frac{n!}{(n+r)!} U_{n+r}(\alpha), \label{eq:I-entry}\\
M_{(\lambda_1,i),(\lambda_2,j)}(\alpha)
&= \sum_r w_{\lambda_1\sqcup\lambda_2,\, r} \frac{n!}{(n+r)!}
\bigl(\Harm_n U_{n+r}(\alpha) - L_{n+r}(\alpha)\bigr)
- \Harm_j I_{(\lambda_1,i),(\lambda_2,j)}(\alpha) \nonumber\\
&\quad + \sum_{\emptyset \neq S \subseteq [\ell(\lambda_2)]}
\frac{j!\,(s-1)!}{(j+s)!}
\sum_r w_{\lambda_1\sqcup\lambda_{2\setminus S},\, r}
\frac{(n+s)!}{(n+s+r)!} U_{n+s+r}(\alpha), \label{eq:M-entry}
\end{align}
with $s := s_{\lambda_2}(S)$.
\end{enumerate}
\end{proposition}

\begin{proof}
Differentiating $\frac{1}{i}\frac{\dd}{\dd\tau}(e^{\Phi_\alpha(z)}/z^{n+1})
= (\gamma+\alpha+1+\log z)e^{\Phi_\alpha(z)}/z^{n+1} - (n+1)e^{\Phi_\alpha(z)}/z^{n+2}$ ($z=1+i\tau$)
and integrating over $[-R, R]$ drops boundary terms $O(R^{-n} e^{-\pi R/2}) \to 0$ by~\eqref{eq:Phi-abs};
conjugation symmetry $\Phi_\alpha(\bar{z}) = \overline{\Phi_\alpha(z)}$ gives~(i).

For~(ii), analytic continuation of Gamma integrals ($\re z = 1$) gives
\begin{equation}\label{eq:gamma-moments}
\int_0^\infty x^n e^{-zx}\dd x = \frac{n!}{z^{n+1}},
\qquad
\int_0^\infty x^n (\log x) e^{-zx}\dd x
= \frac{n!}{z^{n+1}}\bigl(\Harm_n - \gamma - \log z\bigr),
\end{equation}
and substituting $u = x v$ gives the exact polynomial identities ($j \ge 0, s \ge 2$)
\begin{align}
(\cH x^j) &= x^j \int_0^1 \frac{1 - (1-v)^j}{v}\dd v
= x^j \int_0^1 \sum_{r=0}^{j-1}(1-v)^r\dd v = \Harm_j x^j, \nonumber\\
(B_{s-1} x^j) &= x^{j+s}\int_0^1 (1-v)^j v^{s-1}\dd v
= \frac{j!(s-1)!}{(j+s)!} x^{j+s}. \label{eq:H-Bs-monomial}
\end{align}
Inserting $\widetilde{Q}_{\lambda_1}, \widetilde{Q}_{\lambda_2}$, \eqref{eq:WPPP-def}, \eqref{eq:gamma-moments},
and~\eqref{eq:H-Bs-monomial} into~\eqref{eq:I-inf-def}--\eqref{eq:M-inf-def}, and re-indexing
$(\mu_1,\mu_2,\lambda,S)$ with $\lambda_{\setminus S} = \mu_2$ as $(\lambda_1,\lambda_2,S)$ with
$\mu_2 = \lambda_{2\setminus S}$, produces~\eqref{eq:I-entry}--\eqref{eq:M-entry} upon using
$\frac{1}{2\pi}\int_{-\infty}^\infty e^{\Phi_\alpha(z)} z^{-(m+1)}\dd\tau = \frac{1}{m!}U_m(\alpha)$
and symmetrizing $2 c^\top M(\alpha) c = c^\top(M(\alpha)+M(\alpha)^\top)c$.
\end{proof}

\section{Rigorous Arb Certification and Verification}\label{sec:cert}

By Propositions~\ref{prop:contour-limit} and~\ref{prop:UL-reduction}, Theorem~\ref{thm:main}
reduces rigorously to verifying $I(\alpha) \succ 0$ (certifying $\mathcal{I}_\infty(\alpha) = c^\top I(\alpha) c > 0$)
and the Rayleigh lower bound
\begin{equation}\label{eq:rq-cert}
\frac{c^\top D_{\mathrm{net}}(\alpha) c}{c^\top I(\alpha) c} \ge -c^*
\end{equation}
at $\alpha = -523/2000$ for the exact rational witness $c \in \mathbb{Q}^{250}$ (serialized to $60$
significant digits and parsed via exact \texttt{fmpq} rationals before conversion to Arb balls).
While Lemma~\ref{lem:nu-pos} establishes $\mathcal{I}_\infty(\alpha) > 0$ for the radial class
$\Lambda_0$ alone, for the mixed $\Lambda_9$ ansatz positivity (and in fact strict positive
definiteness of the Gram matrix $I(\alpha)$) is rigorously established by interval $LDL^\top$
factorization.

To eliminate single-implementation error, generation/certification (\texttt{cert\_asymptotic.py})
and independent verification (\texttt{verify\_asymptotic.py}) share no code and evaluate the
moments independently:
\begin{itemize}
\item \textbf{Certifier (\texttt{cert\_asymptotic.py}):} Evaluates $U_0(\alpha),\dots,U_{41}(\alpha)$
at $768$-bit precision on $[0, 200]$ (panels $[0,1,4,12,35,90,200]$, adaptive \texttt{acb.integral}
\cite{JohanssonArb}, tolerance $10^{-145}$), adds the tail bound on $[200,\infty)$ from Lemma~\ref{lem:arb-bounds}(i),
and computes $D_{\mathrm{net}}(\alpha)$ via Proposition~\ref{prop:UL-reduction}(i).
\item \textbf{Independent verifier (\texttt{verify\_asymptotic.py}):} Evaluates $U_0(\alpha),\dots,U_{40}(\alpha)$
and \emph{directly} evaluates $L_0(\alpha),\dots,L_{40}(\alpha)$ \emph{without} the IBP identity on
the independent geometric grid $[0, 0.35, 0.9, 2.0, 4.5, 10, 22, 45, 85, 145, 225]$ ($T = 225$) using
$240$-node Gauss--Legendre quadrature at $768$-bit Arb precision. Every node and weight $(x_i, w_i)$
($i = 0, \dots, 239$) is computed inside Arb via \texttt{arb.legendre\_p\_root(240, i, weight=True)},
returning certified midpoint-radius balls with zero external floating-point tables.
\end{itemize}
Numerically, across the $W \le 9$ slice ($30$ blocks, $N = 250$), the asymptotic Rayleigh quotient
varies by only $\approx 10^{-7}$ over $\alpha \in [-0.262, -0.260]$ ($\eta \approx -0.26098864$ at
$\alpha = -0.262$, $-0.2609886366$ at $-0.2615$, and $-0.26098873$ at $-0.260$), making
$\alpha = -523/2000 = -0.2615$ a natural low-denominator rational shift capturing $7$ significant
digits of the flat optimum while keeping concatenated partition lengths $\le 8$ and the $I(\alpha)$
Cholesky pivot ratio $7.90 \times 10^{-11}$.

\begin{lemma}[Analytic Tail and Bernstein-Ellipse Error Bounds]\label{lem:arb-bounds}
Fix $\alpha \in \R$, $T \ge 10$, and integer $n \ge 0$, and put $C_T := \gamma + \frac{\pi}{2} + \log(1+T)$.
\begin{enumerate}
\item[(i)] The truncated tails in~\eqref{eq:Un-Ln-def} on $[T,\infty)$ satisfy
\begin{align}
\left| U_n(\alpha) - \frac{n!}{\pi}\re \int_0^T \frac{e^{\Phi_\alpha(1+i\tau)}}{(1+i\tau)^{n+1}}\dd\tau \right|
&\le \frac{2 e^{\gamma+\alpha+1} n!\, e^{-\pi T/2}}{\pi^2 T^n}, \label{eq:Un-tail}\\
\left| L_n(\alpha) - \frac{n!}{\pi}\re \int_0^T \bigl(\gamma+\log(1+i\tau)\bigr)
\frac{e^{\Phi_\alpha(1+i\tau)}}{(1+i\tau)^{n+1}}\dd\tau \right|
&\le \frac{n!\, e^{\gamma+\alpha+1} e^{-\pi T/2}}{\pi T^n}\left(\frac{2 C_T}{\pi} + \frac{4}{\pi^2 T}\right). \label{eq:Ln-tail}
\end{align}
\item[(ii)] Let $[a,b] = [\mu - h, \mu + h] \subset [0, T]$ ($h > 0$), let $E_\rho$ ($\rho > 1$) be the
Bernstein ellipse with semi-axes $K := \frac{1}{2}(\rho+\rho^{-1})$ and $B := \frac{1}{2}(\rho-\rho^{-1})$,
and let $\overline{D}_\rho$ be its closed interior. If $(1 - h B > 0) \lor (\mu - h K > 0)$, then for
all $u \in \overline{D}_\rho$, $z(u) := 1 + i(\mu + h u)$ avoids the branch cut $(-\infty, 0]$.
\item[(iii)] Under the condition of~(ii), if $f$ is holomorphic on a neighborhood of $\mu + h\overline{D}_\rho$
and $M := \sup_{u \in E_\rho} |f(\mu + h u)|$, then $m$-point Gauss--Legendre quadrature $Q_m$ on $[a,b]$ satisfies
\begin{equation}\label{eq:GL-bernstein-err}
\left|\int_a^b f(t)\dd t - Q_m[f]\right| \le \frac{8 h M}{(\rho - 1)\rho^{2m-1}}.
\end{equation}
\end{enumerate}
\end{lemma}

\begin{proof}
For~(i), $\arctan\tau \ge \frac{\pi}{2} - \frac{1}{\tau}$ and~\eqref{eq:Phi-abs} yield
$|e^{\Phi_\alpha(1+i\tau)}/(1+i\tau)^{n+1}| \le e^{\gamma+\alpha+1} T^{-n} e^{-\pi\tau/2}$ ($\tau \ge T$),
giving~\eqref{eq:Un-tail}. Also $|\gamma+\log(1+i\tau)| \le C_T + \frac{\tau-T}{T}$, whose integral against
$e^{-\pi\tau/2}$ proves~\eqref{eq:Ln-tail}.

For~(ii), any $u = x_u + i y_u \in \overline{D}_\rho$ has $|x_u| \le K$ and $|y_u| \le B$, so
$1 - hB > 0 \implies \re z(u) \ge 1 - hB > 0$ and $\mu - hK > 0 \implies \im z(u) \ge \mu - hK > 0$.

For~(iii), $g(u) := f(\mu+h u)$ pulls back under $u = \frac{1}{2}(w+w^{-1})$ to a holomorphic function
on an open neighborhood of $\{\rho^{-1} \le |w| \le \rho\}$. In its Chebyshev expansion
$g(u) = \sum_{j=0}^\infty a_j T_j(u)$ ($T_j(\frac{w+w^{-1}}{2}) = \frac{1}{2}(w^j+w^{-j})$), matching
Laurent coefficients on $|w| = 1$ gives $c_j = a_j/2$ ($j \ge 1$); shifting the contour to $|w| = \rho$ yields
$a_j = \frac{1}{\pi i}\oint_{|w|=\rho} g(\frac{w+w^{-1}}{2}) w^{-j-1}\dd w$, hence $|a_j| \le 2M\rho^{-j}$ \cite{Trefethen}.
Since $Q_m$ is exact up to degree $2m-1$ and $|\int_{-1}^1 T_j - Q_m[-1,1][T_j]| \le 4$, rescaling gives~\eqref{eq:GL-bernstein-err}.
\end{proof}

Executing \texttt{verify\_asymptotic.py} checks $(1-hB > 0)\lor(\mu-hK > 0)$ on every panel, encloses
$M = \sup_{u \in E_\rho} |G(1+i(\mu+hu))|$ rigorously by partitioning $[0, 2\pi]$ into $32$ angular
intervals $\theta \in [2\pi j/32,\, 2\pi(j+1)/32]$, evaluating each ellipse arc $u(\theta) =
K\cos\theta + i B\sin\theta$ in an Arb complex box \texttt{acb} (checking holomorphicity on every box;
since the holomorphic integrands $G(z)$ in~\eqref{eq:Un-Ln-def} satisfy $\overline{G(\bar{z})} = G(z)$
and $E_\rho = \overline{E_\rho}$, $t \mapsto \frac{1}{2}(G(1+it)+G(1-it))$ extends $\re G(1+it)$
holomorphically to $\mu + h\overline{D}_\rho$ with modulus $\le M$), adds~\eqref{eq:GL-bernstein-err} and~\eqref{eq:Un-tail}--\eqref{eq:Ln-tail},
confirms agreement with all stored moments $U_0,\dots,U_{40}$ (the independently computed
$U_0,\dots,U_{40}$ and $L_0,\dots,L_{40}$ achieve maximum relative Arb radius $1.30 \times 10^{-148}$),
verifies $I(\alpha) \succ 0$ via interval $LDL^\top$, passes all $9/9$ adversarial tamper-rejection tests,
certifies~\eqref{eq:rq-cert}, and establishes
\begin{equation}\label{eq:eta-cert}
\begin{aligned}
\eta_* \in \bigl[&{-0.26098863656469156541919317124794770922900940468892},\\
&{-0.26098863656469156541919317124794770922900940468891}\bigr]
\end{aligned}
\end{equation}
(outer-rounded to $50$ decimal digits from the $768$-bit Arb ball of radius $4.03 \times 10^{-149}$),
proving Theorem~\ref{thm:main}.

As independent cross-checks (non-certified supporting evidence for Proposition~\ref{prop:exact-laplace} and
Proposition~\ref{prop:contour-limit}): comparing Proposition~\ref{prop:exact-laplace} against direct 5D real-space
Duffy simplex quadrature ($14$-point Gauss--Legendre per axis vs.\ $320$-bit Arb Laplace) at $k = 5$ on all $30$
$\Lambda_9$ blocks yields relative errors $5.64 \times 10^{-8}$ ($I$), $8.01 \times 10^{-8}$ ($kJ$), and
$2.37 \times 10^{-8}$ ($R_5(F) = 1.95890046\dots$). Table~\ref{tab:finite-k} evaluates $R_k(F) - \log k$ for
the $\Lambda_9$ witness across $k = 10^3,\dots,10^8$ in $448$-bit Arb Laplace arithmetic ($64$-node composite
Gauss--Legendre, spot-checked against $80$ nodes at $k = 10^6$ with $448$-bit Arb difference $3.71 \times 10^{-23}$),
confirming monotonic shrinkage of the $O(L^r/k)$ remainder $\Delta_L(k) := (R_k(F) - L) - \eta_{L\text{-ref}}(k)$,
where
\[
\eta_{L\text{-ref}}(k) := \frac{L^2\mathcal{I}_\infty(\alpha) + 2L\mathcal{M}_\infty(\alpha) + \mathcal{N}_\infty(\alpha)}{(L+\alpha)\mathcal{I}_\infty(\alpha)} - L
\]
retains the $O(1)$ inner integral $\mathcal{N}_\infty(\alpha) := \frac{1}{2\pi}\int_{-\infty}^\infty e^{\Phi_\alpha(1+i\tau)}
\sum_{\mu_1,\mu_2\in\Lambda_9} W_{\mu_1\sqcup\mu_2}^{\mathrm{PPP}}(1+i\tau)\, N_{\mu_1,\mu_2}^{(2)}(\tau)\dd\tau$,
evaluated in $384$-bit Arb via the second-log moments
$L_n^{(2)}(\alpha) := \frac{n!}{\pi}\re\int_0^\infty (\gamma+\log(1+i\tau))^2 \frac{e^{\Phi_\alpha(1+i\tau)}}{(1+i\tau)^{n+1}}\dd\tau$
and the exact IBP identity $L_n^{(2)}(\alpha) = L_{n+1}(\alpha) - (1+\alpha)L_n(\alpha) - (n+1)^{-1}U_{n+1}(\alpha)$.

\begin{table}[ht]
\centering
\caption{Finite-$k$ convergence of the $\Lambda_9$ witness ($448$-bit Arb Laplace) toward its
$768$-bit certified asymptotic limit, together with the $384$-bit remainder $\Delta_L(k)$ and a
quadratic least-squares fit in $1/(L+\alpha)$ over $k \in \{10^5, 10^6, 10^7, 10^8\}$.}\label{tab:finite-k}
\small
\setlength{\tabcolsep}{7pt}
\begin{tabular}{@{}rccc@{}}
\toprule
$k$ & $R_k(F) - \log k$ & $\Delta_L(k)$ & Gain vs.\ \cite{Lewulis} limit ($-0.273572$) \\
\midrule
$10^3$ & $-0.135735$ & $+1.67 \times 10^{-3}$ & $+0.137837$ \\
$10^4$ & $-0.168933$ & $+2.71 \times 10^{-4}$ & $+0.104639$ \\
$10^5$ & $-0.187944$ & $+4.38 \times 10^{-5}$ & $+0.085628$ \\
$10^6$ & $-0.200383$ & $+6.56 \times 10^{-6}$ & $+0.073189$ \\
$10^7$ & $-0.209188$ & $+9.23 \times 10^{-7}$ & $+0.064384$ \\
$10^8$ & $-0.215757$ & $+1.24 \times 10^{-7}$ & $+0.057815$ \\
\midrule
$1/(L+\alpha)$ fit & $-0.260771$ & --- & $+0.012801$ \\
Certified limit & $\mathbf{-0.260989}$ & --- & $\mathbf{+0.012583}$ \\
\bottomrule
\end{tabular}
\end{table}

\begin{remark}[Open Question]
Combining Theorem~\ref{thm:main} with the barrier~\eqref{eq:pm8b-barrier} gives
\[
0 \le \liminf_{k\to\infty}(\log k - M_k) \le \limsup_{k\to\infty}(\log k - M_k) \le 0.26098863656469156542.
\]
It remains open whether $\lim_{k\to\infty}(\log k - M_k)$ exists and is strictly positive, or
whether $M_k = \log k - o(1)$.
\end{remark}

\section{Reproducibility and Artifact Verification}\label{sec:repro}

The certificate, independent verifier, finite-$k$ checks, and both language versions of this manuscript
are archived at:
\begin{center}
\url{https://doi.org/10.5281/zenodo.24528079}
\end{center}
From the repository root, executing the verification command
\begin{samepage}
\begin{verbatim}
python verify_asymptotic.py cert_asymptotic.json --entry C
\end{verbatim}
\end{samepage}
runs the independent Bernstein-ellipse verifier in approximately $24.5\text{ s}$ (tested on
\texttt{Windows~11}, Python \texttt{3.13.17}, \texttt{python-flint 0.9.0}, peak RSS $77.8\text{ MB}$),
while \texttt{python verify\_all.py} runs the complete suite in $18\text{--}36\text{ s}$ on at most
$4$ CPU threads (with \texttt{numpy 2.5.4}, \texttt{scipy 1.18.1}, and \texttt{mpmath 1.4.1} used
for supplementary checks). Canonical SHA-256 checksums appear in Table~\ref{tab:sha256}.

\begin{table}[ht]
\centering
\caption{Canonical SHA-256 digests of the asymptotic certificate and independent Arb verifier.}\label{tab:sha256}
\small
\begin{tabular}{@{}ll@{}}
\toprule
Artifact & SHA-256 digest \\
\midrule
\texttt{cert\_asymptotic.json} & \texttt{\scriptsize 541e5d4fd45aff5905b38aa30dd5b77cc4d9f11a360823946498f26e80e02921} \\
\texttt{verify\_asymptotic.py} & \texttt{\scriptsize e3721c23fdeaec99c95d9a6d34ae83ee9516410f6ce035d6cf53efaf53970a78} \\
\bottomrule
\end{tabular}
\end{table}

\subsection*{Computational Assistance and Verification}
Numerical exploration and code prototyping were assisted by computational and AI tools. Every proof
in this manuscript was audited analytically, and every reported asymptotic constant and moment bound
is independently verified deterministically via Arb ball arithmetic (\texttt{verify\_asymptotic.py}).

\begin{thebibliography}{9}\raggedright

\bibitem{GPY}
D.~A.~Goldston, J.~Pintz, and C.~Y.~Y{\i}ld{\i}r{\i}m,
\emph{Primes in tuples.~I},
Ann. of Math. (2) \textbf{170} (2009), no.~2, 819--862.

\bibitem{JohanssonArb}
F.~Johansson,
\emph{Arb: efficient arbitrary-precision midpoint-radius interval arithmetic},
IEEE Trans. Comput. \textbf{66} (2017), no.~8, 1281--1292.

\bibitem{Lewulis}
P.~Lewulis,
\emph{Clusters of primes revisited},
preprint (October 2026), \href{https://arxiv.org/abs/2610.41113}{arXiv:2610.41113}.

\bibitem{Maynard}
J.~Maynard,
\emph{Small gaps between primes},
Ann. of Math. (2) \textbf{181} (2015), no.~1, 383--413.

\bibitem{PM8b}
D.~H.~J.~Polymath,
\emph{Variants of the Selberg sieve, and bounded intervals containing many primes},
Res. Math. Sci. \textbf{1} (2014), Art.~12, 83 pp.

\bibitem{Stadlmann2025}
J.~Stadlmann,
\emph{On primes in arithmetic progressions and bounded gaps between many primes},
Adv. Math. \textbf{468} (2025), Art.~110190, \href{https://doi.org/10.1016/j.aim.2025.110190}{doi:10.1016/j.aim.2025.110190}
(\href{https://arxiv.org/abs/2309.00425}{arXiv:2309.00425}).

\bibitem{Stadlmann2026}
J.~Stadlmann,
\emph{Bounded gaps between primes},
preprint (August 2026), \href{https://arxiv.org/abs/2608.31126}{arXiv:2608.31126}.

\bibitem{Trefethen}
L.~N.~Trefethen,
\emph{Is Gauss quadrature better than Clenshaw--Curtis?},
SIAM Rev. \textbf{50} (2008), no.~1, 67--87.

\bibitem{Zhang}
Y.~Zhang,
\emph{Bounded gaps between primes},
Ann. of Math. (2) \textbf{179} (2014), no.~3, 1121--1174.

\end{thebibliography}

\end{document}
